\chapter{Pràctica}
\label{cap:practica}

En aquest capítol es recullen exercicis resolts que permeten consolidar els conceptes principals de la teoria. Cada exercici presenta primer l'enunciat i, a continuació, la solució. L'enunciat no conté cap entorn de solució.

\section{Definicions. Anell dels polinomis}

\begin{exercici}
Indica quines de les expressions següents són monomis i, en cas
afirmatiu, indical seu grau i coeficient.\[
\begin{array}{lllllll}
a)~-4x^{3} & & b)~-\frac{25}{x^{3}} & & c)~-\frac{2y^{4}}{3} & & d)~-\frac{1}{2} \\
& & & & & & \\
e)~3z^{3}-2z^{2} & & f)~4\sqrt[3]{x^{2}} & & g)~4\sqrt{2}t~ & &
h)~-2x^{2}+x^{2}\end{array}\]
\end{exercici}

\begin{solucio}
\begin{enumerate}
\item[a)] Es tracta d'un monomi en la indeterminada $x$ de grau $3$ i
coeficient $-4$.

\item[b)] No es tracta d'un monomi perquè\[
-\frac{25}{x^{3}}=-25x^{-3}
\]i l'exponent de $x$ no és un nombre natural.

\item[c)] Es tracta d'un monomi en la indeterminada $y$ de grau $4$ i
coeficient $-\frac{2}{3}$, ja que\[
-\frac{2y^{4}}{3}=-\frac{2}{3}y^{4}
\]

\item[d)] Es tracta d'un monomi de grau zero i coeficient $-\frac{1}{2}$
en una certa indeterminada.

\item[e)] Per la forma de l'expressió, no es tracta d'un monomi; en
realitat és un binomi, és a dir, la suma de dos monomis, $3z^{3}$ i $-2z^{2}.$

\item[f)] No es tracta d'un monomi perquè\[
4\sqrt[3]{x^{2}}=4x^{\frac{2}{3}}
\]i l'exponent de $x$ no és un nombre natural.

\item[g)] Es tracta d'un monomi en la indeterminada $t$ de grau $1$ i
coeficient $4\sqrt{2}$.

\item[h)] Es tracta d'un monomi en la indeterminada $x$ de grau $2$ i i
coeficient $-1$, perquè\[
-2x^{2}+x^{2}=(-2+1)x^{2}=-x^{2}
\]
\end{enumerate}
\end{solucio}

\begin{exercici}
Indica quines de les expressions següents són monomis en la
indeterminada $x$, i, en cas afirmatiu, indica'n el grau i coeficient.\[
\begin{array}{lllllll}
a)~3x^{3}\cdot x^{2} & & b)~3x^{4}+x^{2} & & c)~-\frac{3x^{5}}{x^{3}} & &
d)~\sqrt[3]{-\frac{x^{9}}{27}} \\
& & & & & & \\
e)~3x^{5}-\frac{x^{5}}{2} & & f)~-\frac{4x^{3}}{x^{4}} & &
g)~(x^{2}+x^{3})^{2}~ & & h)~(3x^{2}+x^{2})^{3}\end{array}\]
\end{exercici}

\begin{solucio}

\begin{enumerate}
\item[a)] Es tracta d'un monomi de grau $5$ i coeficient $3$, ja que\[
3x^{3}\cdot x^{2}=3x^{5}
\]

\item[b)] Per la forma de l'expressió, no es tracta d'un monomi; en
realitat és un binomi de grau $4$.

\item[c)] Es tracta d'un monomi de grau $2$ i coeficient $-3$, perquè\[
-\frac{3x^{5}}{x^{3}}=-3x^{2}
\]

\item[d)] Es tracta d'un monomi de grau $3$ i coeficient $-\frac{1}{3}$,
perquè\[
\sqrt[3]{-\frac{x^{9}}{27}}=\sqrt[3]{\left( -\frac{x^{3}}{3}\right) ^{3}}=-\frac{1}{3}x^{3}
\]

\item[e)] Es tracta d'un monomi de grau $5$ i coeficient $\frac{5}{2}$,
perquè\[
3x^{5}-\frac{x^{5}}{2}=\left( 3-\frac{1}{2}\right) x^{5}=\frac{5}{2}x^{5}
\]

\item[f)] No es tracta d'un monomi ja que\[
-\frac{4x^{3}}{x^{4}}=-4x^{3-4}=-4x^{-1}
\]és a dir, l'exponent de $x$ no és un nombre natural.

\item[g)] No es tracta d'un monomi, perquè\begin{eqnarray*}
(x^{2}+x^{3})^{2} &=&(x^{2})^{2}+2x^{2}\cdot x^{3}+(x^{3})^{2} \\
&=&x^{4}+2x^{5}+x^{6}
\end{eqnarray*}és a dir, és un trinomi de grau $6$.

\item[h)] Es tracta d'un monomi de grau $6$ i coeficient $64$.\[
(3x^{2}+x^{2})^{3}=(4x^{2})^{3}=64x^{6}
\]
\end{enumerate}
\end{solucio}

\begin{exercici}
Indica quines de les expressions següents són polinomis en la
indeterminada $x$ i, en cas afirmatiu, quin és el seu grau, el seu
coeficient de grau $3$ i el seu terme independent:\[
\begin{array}{llll}
a)~-3x^{3}+4x^{5}-2x+\frac{1}{3} & & & b)~-x^{3}+x^{4}-\frac{1}{x}+1 \\
& & & \\
c)~\sqrt{3}x-4x^{5}+\frac{3x}{11} & & & d)~4x^{5}-\sqrt{3x}+2~\end{array}\]
\end{exercici}

\begin{solucio}

\begin{enumerate}
\item[a)] Es tracta d'un polinomi de grau $5$, perquè\[
\frac{1}{3}-2x-3x^{3}+4x^{5}
\]A més, el seu coeficient de grau $3$ és $-3$ i el seu terme
independent és $\frac{1}{3}$.

\item[b)] No es tracta d'un polinomi, perquè, el terme
\[
\frac{1}{x}=x^{-1}
\]no és un monomi.

\item[c)] És un polinomi de grau $5$, ja que\[
\sqrt{3}x-4x^{5}+\frac{3x}{11}=\left( \sqrt{3}+\frac{3}{11}\right) x-4x^{5}
\]És clar que el coeficient de grau $3$ és $0$ i el seu terme
independent és també $0$.

\item[d)] No es tracta d'un polinomi, perquè el terme
\[
-\sqrt{3x}=-\sqrt{3}x^{\frac{1}{2}}
\]no és un monomi.
\end{enumerate}
\end{solucio}

\begin{exercici}
Per a cadascun dels polinomis següents indica quin dels anells de
polinomis $\mathbb{Z}[x],\mathbb{Q}[x]$ o $\mathbb{R}[x]$ és el més
petit que el conté:\[
\begin{array}{lll}
a)~1-x+x^{2}+x^{4} & & b)~1+x+\sqrt{3}\cdot x^{2} \\
c)~x-x^{2} & & d)~\frac{1}{3}x+4.25x^{10} \\
e)~\sqrt{2}+\frac{3}{4}x-5x^{2} & &
\end{array}\]
\end{exercici}

\begin{solucio}
a) Com que tots els coeficients del polinomi són nombres enters,
es té\[
1-x+x^{2}+x^{4}\in \mathbb{Z}[x]
\]

b) Com que és el coeficient del monomi de grau més gran un nombre
irracional, es té\[
1+x+\sqrt{3}\cdot x^{2}\in \mathbb{R}[x]
\]

c) Com que tots els coeficients del polinomi són nombres enters,
es té\[
x-x^{2}\in \mathbb{Z}[x]
\]

d) Observa que\[
\frac{1}{3}x+4.25x^{10}=\frac{1}{3}x+\frac{17}{4}x^{10}
\]i com que són tots els coeficients del polinomi nombres racionals, es
té\[
\frac{1}{3}x+4.25x^{10}\in \mathbb{Q}[x]
\]

e) Com que és el terme independent un nombre irracional, es té\[
\sqrt{2}+\frac{3}{4}x-5x^{2}\in \mathbb{R}[x]
\]
\end{solucio}

\begin{exercici}
Redueix termes semblants i ordena en potències decreixents les
següents expressions:\[
\begin{array}{l}
a)~3x-x^{3}+4x+2x^{3}-x^{2}+2x-7x^{2}+4x^{3}+8 \\
b)~4x^{5}-2x^{3}+3x^{2}+\frac{1}{2}x^{3}-x^{5}+x^{2}-3\end{array}\]
\end{exercici}

\begin{solucio}
a)
\begin{eqnarray*}
&&3x-x^{3}+4x+2x^{3}-x^{2}+2x-7x^{2}+4x^{3}+8 \\
&=&8+(3+4+2)x-(1+7)x^{2}+(-1+2+4)x^{3} \\
&=&8+9x-8x^{2}+5x^{3}
\end{eqnarray*}

b)
\begin{eqnarray*}
&&4x^{5}-2x^{3}+3x^{2}+\frac{1}{2}x^{3}-x^{5}+x^{2}-3 \\
&=&-3+(3+1)x^{2}+\left( -2+\frac{1}{2}\right) x^{3}+(4-1)x^{5} \\
&=&-3+4x^{2}-\frac{3}{2}x^{3}+3x^{5}
\end{eqnarray*}
\end{solucio}

\begin{exercici}
Troba $m$ i $n$ perquè els polinomis següents siguin iguals\[
\begin{array}{l}
p(x)=(3+m)\cdot x^{2}+7x+21 \\
q(x)=7x^{2}+(2+n)\cdot x+21\end{array}\]
\end{exercici}

\begin{solucio}
Sabem que dos polinomis són iguals si coincideixen els coeficients de
tots els termes del mateix grau en un i l'altre polinomi. Així,
tenim\begin{eqnarray*}
3+m &=&7 \\
7 &=&2+n
\end{eqnarray*}d'on $m=4$ i $n=5$.
\end{solucio}

\begin{exercici}
 Existeixen valors de $A,B$ i $C$ perquè es compleixi la
següent igualtat?\[
3x^{2}-8x+10=Ax^{2}+Bx+A-Cx+Bx^{2}+C
\]
\end{exercici}

\begin{solucio}
Agrupant termes semblants en un membre de la igualtat, obtenim\[
3x^{2}-8x+10=(A+B)x^{2}+(B-C)x+A+C
\]Com que es tracta d’una igualtat de polinomis s'han de complir les següents
equacions simultàniament\begin{eqnarray*}
A+B &=&3 \\
B-C &=&-8 \\
A+C &=&10
\end{eqnarray*}Raquestando les dues primeras, obtenim\[
A+C=11
\]cosa que no és possible, ja que $A+C=10$. Per consegüent, no hi ha valors de $A,B$ i $C$ perquè es compleix la igualtat.
\end{solucio}

\begin{exercici}
Troba el valor numèric dels polinomis següents per als $x$
indicats:\[
\begin{array}{llll}
a) & p(x)=2+\frac{3}{4}x-x^{2} & & x=-1/3 \\
b) & q(x)=\sqrt{3}+\sqrt{3}x^{3} & & x=\sqrt{3} \\
c) & r(x)=\frac{1}{3}x+\frac{1}{2}x^{2}+\frac{4}{5}x^{3} & & x=-\sqrt{2}\end{array}\]
\end{exercici}

\begin{solucio}
a)
\begin{eqnarray*}
p(-1/3) &=&2+\frac{3}{4}\cdot \left( -\frac{1}{3}\right) -\left( -\frac{1}{3}\right) ^{2} \\
&=&2-\frac{1}{4}-\frac{1}{9} \\
&=&\frac{59}{36}
\end{eqnarray*}

b)\begin{eqnarray*}
q(\sqrt{3}) &=&\sqrt{3}+\sqrt{3}\left( \sqrt{3}\right) ^{3} \\
&=&\sqrt{3}+\left( \sqrt{3}\right) ^{4} \\
&=&\sqrt{3}+9
\end{eqnarray*}

c)\begin{eqnarray*}
r(-\sqrt{2}) &=&\frac{1}{3}\left( -\sqrt{2}\right) +\frac{1}{2}\left( -\sqrt{2}\right) ^{2}+\frac{4}{5}\left( -\sqrt{2}\right) ^{3} \\
&=&-\frac{\sqrt{2}}{3}+\frac{1}{2}\cdot 2+\frac{4}{5}\cdot \left( -2\sqrt{2}\right) \\
&=&-\frac{\sqrt{2}}{3}+1-\frac{8\sqrt{2}}{5} \\
&=&1-\frac{29\sqrt{2}}{15}
\end{eqnarray*}
\end{solucio}

\begin{exercici}
Calcula $m$ perquè el valor numèric de $p(x)=3x^{6}+7x^{4}+mx^{2}+2mx+1$ sigui $121$ per a $x=2$.
\end{exercici}

\begin{solucio}
La condició imposada per l’enunciat és equivalent a escriure\[
p(2)=121
\]Aleshores,\begin{eqnarray*}
3\cdot 2^{6}+7\cdot 2^{4}+m\cdot 2^{2}+2m\cdot 2+1 &=&121 \\
192+112+4m+4m+1 &=&121 \\
8m &=&-184 \\
m &=&-23
\end{eqnarray*}
\end{solucio}

\begin{exercici}
Donats els polinomis\[
\begin{array}{l}
p(x)=3-x+\frac{1}{2}x^{2} \\
q(x)=4+5x^{3}+6x^{4} \\
r(x)=3-x^{4}\end{array}\]calcula\[
\begin{array}{l}
a)~p(x)+q(x)-r(x) \\
b)~p(x)-q(x)+r(x) \\
c)~-p(x)+q(x)+r(x)\end{array}\]
\end{exercici}

\begin{solucio}
a)\[
\begin{array}{rccccccccccc}
p(x) & = & & 3 & - & x & + & \frac{1}{2}x^{2} & & & & \\
q(x) & = & & 4 & & & & & + & 5x^{3} & + & 6x^{4} \\
-r(x) & = & - & 3 & & & & & & & + & x^{4} \\ \hline
p(x)+q(x)-r(x) & = & & 4 & - & x & + & \frac{1}{2}x^{2} & + & 5x^{3} & + &
7x^{4}\end{array}\]

b)\[
\begin{array}{rccccccccccc}
p(x) & = & & 3 & - & x & + & \frac{1}{2}x^{2} & & & & \\
-q(x) & = & - & 4 & & & & & - & 5x^{3} & - & 6x^{4} \\
r(x) & = & & 3 & & & & & & & - & x^{4} \\ \hline
p(x)-q(x)+r(x) & = & & 2 & - & x & + & \frac{1}{2}x^{2} & - & 5x^{3} & - &
7x^{4}\end{array}\]

c)\[
\begin{array}{rccccccccccc}
-p(x) & = & - & 3 & + & x & - & \frac{1}{2}x^{2} & & & & \\
q(x) & = & & 4 & & & & & + & 5x^{3} & + & 6x^{4} \\
r(x) & = & & 3 & & & & & & & - & x^{4} \\ \hline
-p(x)+q(x)+r(x) & = & & 4 & + & x & - & \frac{1}{2}x^{2} & + & 5x^{3} & + &
5x^{4}\end{array}\]
\end{solucio}

\begin{exercici}
Troba el polinomi $p(x)$ que sumat amb $q(x)=3-x^{2}+5x^{3}$ dona el
polinomi $r(x)=8+x^{2}+5x^{3}+x^{4}$.
\end{exercici}

\begin{solucio}
La condició de l'enunciat es tradueix de la manera següent:\begin{eqnarray*}
p(x)+q(x) &=&r(x) \\
p(x)+\left( 3-x^{2}+5x^{3}\right) &=&8+x^{2}+5x^{3}+x^{4} \\
p(x) &=&8+x^{2}+5x^{3}+x^{4}-\left( 3-x^{2}+5x^{3}\right) \\
p(x) &=&8+x^{2}+5x^{3}+x^{4}-3+x^{2}-5x^{3} \\
p(x) &=&5+2x^{2}+x^{4}
\end{eqnarray*}
\end{solucio}

\begin{exercici}
Donats els polinomis\[
\begin{array}{l}
p(x)=\sqrt{3}+5x-\sqrt{3}x^{2} \\
q(x)=-\sqrt{3}+2\sqrt{3}x \\
r(x)=2\sqrt{3}-5\sqrt{3}x+x^{3}\end{array}\]calcula\[
\begin{array}{l}
a)~-\left[ p(x)+q(x)\right] +r(x) \\
b)~p(x)-\left[ p(x)+q(x)+r(x)\right] \\
c)~q(x)-r(x) \\
d)~p(x)+q(x)-r(x)\end{array}\]
\end{exercici}

\begin{solucio}
a) Com que\[
-\left[ p(x)+q(x)\right] +r(x)=-p(x)-q(x)+r(x)
\]tenim\[
\begin{array}{rccccccccc}
-p(x) & = & - & \sqrt{3} & - & 5x & + & \sqrt{3}x^{2} & & \\
-q(x) & = & & \sqrt{3} & - & 2\sqrt{3}x & & & & \\
r(x) & = & & 2\sqrt{3} & - & 5\sqrt{3}x & & & + & x^{3} \\ \hline
-p(x)-q(x)+r(x) & = & & 2\sqrt{3} & - & (5+7\sqrt{3})x & + & \sqrt{3}x^{2}
& + & x^{3}\end{array}\]

b) Com que\begin{eqnarray*}
p(x)-\left[ p(x)+q(x)+r(x)\right] &=&p(x)-p(x)-q(x)-r(x) \\
&=&-q(x)-r(x)
\end{eqnarray*}tenim\begin{eqnarray*}
-q(x)-r(x) &=&-\left( -\sqrt{3}+2\sqrt{3}x\right) -\left( 2\sqrt{3}-5\sqrt{3}x+x^{3}\right) \\
&=&\sqrt{3}-2\sqrt{3}x-2\sqrt{3}+5\sqrt{3}x-x^{3} \\
&=&-\sqrt{3}+3\sqrt{3}x-x^{3}
\end{eqnarray*}

c)\begin{eqnarray*}
q(x)-r(x) &=&-\sqrt{3}+2\sqrt{3}x-\left( 2\sqrt{3}-5\sqrt{3}x+x^{3}\right) \\
&=&-\sqrt{3}+2\sqrt{3}x-2\sqrt{3}+5\sqrt{3}x-x^{3} \\
&=&-3\sqrt{3}+7\sqrt{3}x-x^{3}
\end{eqnarray*}

d)
\[
\begin{array}{rccccccccc}
p(x) & = & & \sqrt{3} & + & 5x & - & \sqrt{3}x^{2} & & \\
q(x) & = & - & \sqrt{3} & + & 2\sqrt{3}x & & & & \\
-r(x) & = & - & 2\sqrt{3} & + & 5\sqrt{3}x & & & - & x^{3} \\ \hline
p(x)+q(x)-r(x) & = & - & 2\sqrt{3} & + & (5+7\sqrt{3})x & - & \sqrt{3}x^{2}
& - & x^{3}\end{array}\]
\end{solucio}

\begin{exercici}
Efectual producte de les parelles de polinomis indicades:\[
\begin{array}{l}
a)~3+x-x^{2}+\sqrt{3}x^{3}\text{ i }x+3x^{2}-5x^{3} \\
b)~\sqrt{2}x+2\sqrt{2}x^{2}-\sqrt{2}x^{3}\text{ i }\sqrt{2}x+\sqrt{2}x^{2}\end{array}\]
\end{exercici}

\begin{solucio}
a)\[
\begin{array}{ccccccccccc}
& & & & 3 & + & x & - & x^{2} & + & \sqrt{3}x^{3} \\
& & & & & & x & + & 3x^{2} & - & 5x^{3} \\ \hline
3x & + & x^{2} & - & x^{3} & + & \sqrt{3}x^{4} & & & & \\
& & 9x^{2} & + & 3x^{3} & - & 3x^{4} & + & 3\sqrt{3}x^{5} & & \\
& & & - & 15x^{3} & - & 5x^{4} & + & 5x^{5} & - & 5\sqrt{3}x^{6} \\ \hline
3x & + & 10x^{2} & - & 13x^{3} & + & (\sqrt{3}-8)x^{4} & + & (3\sqrt{3}+5)x^{5} & - & 5\sqrt{3}x^{6}\end{array}\]

b) Com que\begin{eqnarray*}
\sqrt{2}x+2\sqrt{2}x^{2}-\sqrt{2}x^{3} &=&\sqrt{2}x\left( 1+2x-x^{2}\right)
\\
\sqrt{2}x+\sqrt{2}x^{2} &=&\sqrt{2}x(1+x)
\end{eqnarray*}\[
\begin{array}{ccccccc}
& & 1 & + & 2x & - & x^{2} \\
& & & & 1 & + & x \\ \hline
1 & + & 2x & - & x^{2} & & \\
& & x & + & 2x^{2} & - & x^{3} \\ \hline
1 & + & 3x & + & x^{2} & - & x^{3}\end{array}\]Aleshores,\begin{eqnarray*}
&&\left( \sqrt{2}x+2\sqrt{2}x^{2}-\sqrt{2}x^{3}\right) \cdot \left( \sqrt{2}x+\sqrt{2}x^{2}\right) \\
&=&\sqrt{2}x\left( 1+2x-x^{2}\right) \cdot \sqrt{2}x(1+x) \\
&=&\left( \sqrt{2}\right) ^{2}x^{2}\cdot \left( 1+3x+x^{2}-x^{3}\right) \\
&=&2x^{2}\cdot \left( 1+3x+x^{2}-x^{3}\right) \\
&=&2x^{2}+6x^{3}+2x^{4}-2x^{5}
\end{eqnarray*}
\end{solucio}

\begin{exercici}
Efectua els productes següents:\[
\begin{array}{l}
a)~(x-1)\cdot (x+1) \\
b)~(x-a)\cdot (x+a) \\
c)~(x+3)\cdot (x+1)\cdot (x-3) \\
d)~(-x-1)\cdot (-x-1)\end{array}\]
\end{exercici}

\begin{solucio}
a) Es tracta d’una suma per diferència i, per tant, tenim\[
(x-1)\cdot (x+1)=x^{2}-1
\]

b) Anàlogament, tenim\[
(x-a)\cdot (x+a)=x^{2}-a^{2}
\]

c) De la mateixa manera, tenim\begin{eqnarray*}
(x+3)\cdot (x+1)\cdot (x-3) &=&(x^{2}-9)\cdot (x+1) \\
&=&x^{3}+x^{2}-9x-9
\end{eqnarray*}

d) Es tracta del quadrat d'un binomi i, per tant, tenim\begin{eqnarray*}
(-x-1)\cdot (-x-1) &=&(-1)^{2}\cdot (x+1)^{2} \\
&=&x^{2}+2x+1
\end{eqnarray*}
\end{solucio}

\begin{exercici}
Donats els polinomis\[
\begin{array}{l}
p(x)=x^{3}-\sqrt{2}x+3 \\
q(x)=2x-7x^{3} \\
r(x)=1+x^{2}\end{array}\]calcula\[
\begin{array}{l}
a)~2\cdot p(x)-q(x)+r(x) \\
b)~p(x)\cdot q(x)-2\cdot p(x) \\
c)~p(x)\cdot p(x)-q(x)\cdot q(x) \\
d)~p(x)\cdot r(x)-q(x)\cdot r(x)\end{array}\]
\end{exercici}

\begin{solucio}
a)\begin{eqnarray*}
2\cdot p(x)-q(x)+r(x) &=&2\left( x^{3}-\sqrt{2}x+3\right) -\left(
2x-7x^{3}\right) +1+x^{2} \\
&=&2x^{3}-2\sqrt{2}x+6-2x+7x^{3}+1+x^{2} \\
&=&(2+7)x^{3}+x^{2}-(2\sqrt{2}+2)x+6+1 \\
&=&9x^{3}+x^{2}-(2\sqrt{2}+2)x+7
\end{eqnarray*}

b)\begin{eqnarray*}
p(x)\cdot q(x)-2\cdot p(x) &=&\left( x^{3}-\sqrt{2}x+3\right) \cdot \left(
2x-7x^{3}\right) -2\left( x^{3}-\sqrt{2}x+3\right) \\
&=&2x^{4}-7x^{6}-2\sqrt{2}x^{2}+7\sqrt{2}x^{4}+6x-21x^{3}-2x^{3}+2\sqrt{2}x-6
\\
&=&-7x^{6}+(2+7\sqrt{2})x^{4}-23x^{3}-2\sqrt{2}x^{2}+(6+2\sqrt{2})x-6
\end{eqnarray*}

c)\begin{eqnarray*}
p(x)\cdot p(x)-q(x)\cdot q(x) &=&\left[ p(x)\right] ^{2}-\left[ q(x)\right]
^{2} \\
&=&\left( x^{3}-\sqrt{2}x+3\right) ^{2}-\left( 2x-7x^{3}\right) ^{2} \\
&=&-48x^{6}+(28-2\sqrt{2})x^{4}+6x^{3}-2x^{2}-6\sqrt{2}x+9
\end{eqnarray*}

d)
\begin{eqnarray*}
p(x)\cdot r(x)-q(x)\cdot r(x) &=&\left[ p(x)-q(x)\right] \cdot r(x) \\
&=&\left( x^{3}-\sqrt{2}x+3-2x+7x^{3}\right) \cdot \left( 1+x^{2}\right) \\
&=&\left[ 8x^{3}-\left( 2+\sqrt{2}\right) x+3\right] \cdot \left(
1+x^{2}\right) \\
&=&8x^{5}+\left( 6-\sqrt{2}\right) x^{3}+3x^{2}-\left( 2+\sqrt{2}\right) x+3
\end{eqnarray*}
\end{solucio}

\begin{exercici}
Efectua les operacions següents:\[
\begin{array}{l}
a)~x\cdot (x^{2}-5)^{2}-(x+1)\cdot (x+3)^{2} \\
b)~\left( \frac{2}{3}x+5\right) \cdot \left( \frac{2}{3}x-5\right) -\left(
\frac{1}{3}x+5\right) ^{2} \\
c)~(3x+5)\cdot (x^{2}-x+3)-(2x+5)\cdot (x^{2}-x+3)\end{array}\]
\end{exercici}

\begin{solucio}
a)\begin{eqnarray*}
x\cdot (x^{2}-5)^{2}-(x+1)\cdot (x+3)^{2} &=&x\cdot \left(
x^{4}-10x^{2}+25\right) -(x+1)\cdot (x^{2}+6x+9) \\
&=&x^{5}-10x^{3}+25x-x^{3}-6x^{2}-9x-x^{2}-6x-9 \\
&=&x^{5}-11x^{3}-7x^{2}+10x-9
\end{eqnarray*}

b)\begin{eqnarray*}
\left( \frac{2}{3}x+5\right) \cdot \left( \frac{2}{3}x-5\right) -\left(
\frac{1}{3}x+5\right) ^{2} &=&\frac{4}{9}\cdot x^{2}-25-\frac{1}{9}\cdot
x^{2}-\frac{10}{3}\cdot x-25 \\
&=&\frac{1}{3}\cdot x^{2}-\frac{10}{3}\cdot x-50
\end{eqnarray*}

c)\begin{eqnarray*}
(3x+5)\cdot (x^{2}-x+3)-(2x+5)\cdot (x^{2}-x+3) &=&(3x+5-2x-5)\cdot
(x^{2}-x+3) \\
&=&x\cdot (x^{2}-x+3) \\
&=&x^{3}-x^{2}+3x
\end{eqnarray*}
\end{solucio}

\section{Potència d'un binomi. Fórmula de Newton}

\begin{exercici}
Calcula\[
\begin{array}{ccccccc}
a)~\left( 3x-1\right) ^{3} & & b)~\left( \sqrt{2}+\frac{3}{x}\right) ^{4} &
& c)~\left( 2\sqrt{2}+3\sqrt{3}\right) ^{4} & & d)~\left( \sqrt{2}x+\frac{x^{2}}{2}\right) ^{3}\end{array}\]
\end{exercici}

\begin{solucio}
a) Mitjançant la fórmula de Newton, obtenim\begin{eqnarray*}
\left( 3x-1\right) ^{3} &=&\binom{3}{0}(3x)^{3}+\binom{3}{1}(3x)^{2}\cdot
(-1)+\binom{3}{2}3x\cdot (-1)^{2}+\binom{3}{3}(-1)^{3} \\
&=&\binom{3}{0}\cdot 27x^{3}-\binom{3}{1}\cdot 9x^{2}+\binom{3}{2}\cdot 3x-\binom{3}{3}
\end{eqnarray*}Ara calcularem els nombres combinatoris mitjançant el triangle
de Tartaglia\[
\begin{array}{ccccccc}
& & 1 & & 1 & & \\
& 1 & & 2 & & 1 & \\
1 & & 3 & & 3 & & 1\end{array}\]Aleshores,\[
\left( 3x-1\right) ^{3}=27x^{3}-27x^{2}+9x-1
\]

b) Mitjançant la fórmula de Newton, obtenim\begin{eqnarray*}
\left( \sqrt{2}+\frac{3}{x}\right) ^{4} &=&\binom{4}{0}(\sqrt{2})^{4}+\binom{4}{1}(\sqrt{2})^{3}\cdot \frac{3}{x}+\binom{4}{2}(\sqrt{2})^{2}\cdot \left(
\frac{3}{x}\right) ^{2}+\binom{4}{3}\sqrt{2}\cdot \left( \frac{3}{x}\right)
^{3}+\binom{4}{4}\left( \frac{3}{x}\right) ^{4} \\
&=&\binom{4}{0}\cdot 4+\binom{4}{1}\cdot \frac{6\sqrt{2}}{x}+\binom{4}{2}\cdot \frac{18}{x^{2}}+\binom{4}{3}\cdot \frac{27\sqrt{2}}{x^{3}}+\binom{4}{4}\cdot \frac{81}{x^{4}}
\end{eqnarray*}Ara calcularem els nombres combinatoris mitjançant el triangle
de Tartaglia\[
\begin{array}{ccccccccc}
& & & 1 & & 1 & & & \\
& & 1 & & 2 & & 1 & & \\
& 1 & & 3 & & 3 & & 1 & \\
1 & & 4 & & 6 & & 4 & & 1\end{array}\]Aleshores,\[
\left( \sqrt{2}+\frac{3}{x}\right) ^{4}=4+\frac{24\sqrt{2}}{x}+\frac{108}{x^{2}}+\frac{108\sqrt{2}}{x^{3}}+\frac{81}{x^{4}}
\]

c) De la mateixa manera, obtenim\begin{eqnarray*}
\left( 2\sqrt{2}+3\sqrt{3}\right) ^{4} &=&\binom{4}{0}(2\sqrt{2})^{4}+\binom{4}{1}(2\sqrt{2})^{3}\cdot 3\sqrt{3}+\binom{4}{2}(2\sqrt{2})^{2}\cdot \left( 3\sqrt{3}\right) ^{2}+\binom{4}{3}2\sqrt{2}\cdot \left( 3\sqrt{3}\right) ^{3}+\binom{4}{4}\left( 3\sqrt{3}\right) ^{4} \\
&=&64+192\sqrt{6}+1296+648\sqrt{6}+729 \\
&=&2089+840\sqrt{6}
\end{eqnarray*}

d) De la mateixa manera, obtenim\begin{eqnarray*}
\left( \sqrt{2}x+\frac{x^{2}}{2}\right) ^{3} &=&\left[ \left( \sqrt{2}+\frac{x}{2}\right) \cdot x\right] ^{3} \\
&=&x^{3}\cdot \left( \sqrt{2}+\frac{x}{2}\right) ^{3} \\
&=&x^{3}\cdot \left[ \binom{3}{0}(\sqrt{2})^{3}+\binom{3}{1}(\sqrt{2})^{2}\cdot \frac{x}{2}+\binom{3}{2}\sqrt{2}\cdot \left( \frac{x}{2}\right)
^{2}+\binom{3}{3}\left( \frac{x}{2}\right) ^{3}\right] \\
&=&x^{3}\cdot \left( 2\sqrt{2}+3x+\frac{3\sqrt{2}}{4}x^{2}+\frac{1}{8}x^{3}\right) \\
&=&2\sqrt{2}\cdot x^{3}+3\cdot x^{4}+\frac{3\sqrt{2}}{4}\cdot x^{5}+\frac{1}{8}\cdot x^{6}
\end{eqnarray*}
\end{solucio}

\begin{exercici}
Calcula\[
\begin{array}{l}
a)~\binom{n}{0}+\binom{n}{1}+\binom{n}{2}+\cdots +\binom{n}{n-1}+\binom{n}{n}
\\
b)~\binom{n}{0}-\binom{n}{1}+\binom{n}{2}+\cdots +(-1)^{n-1}\cdot \binom{n}{n-1}+(-1)^{n}\cdot \binom{n}{n}\end{array}\]
\end{exercici}

\begin{solucio}
a) Observa que es compleix\[
\left( 1+1\right) ^{n}=\binom{n}{0}+\binom{n}{1}+\binom{n}{2}+\cdots +\binom{n}{n-1}+\binom{n}{n}
\]Aleshores,\[
\binom{n}{0}+\binom{n}{1}+\binom{n}{2}+\cdots +\binom{n}{n-1}+\binom{n}{n}=2^{n}
\]

b) De la mateixa manera, tenim\[
\left( 1-1\right) ^{n}=\binom{n}{0}-\binom{n}{1}+\binom{n}{2}+\cdots
+(-1)^{n-1}\cdot \binom{n}{n-1}+(-1)^{n}\cdot \binom{n}{n}
\]Aleshores,\[
\binom{n}{0}-\binom{n}{1}+\binom{n}{2}+\cdots +(-1)^{n-1}\cdot \binom{n}{n-1}+(-1)^{n}\cdot \binom{n}{n}=0
\]
\end{solucio}

\begin{exercici}
Del desenvolupament de $\left( 2x+\frac{1}{2}\right) ^{15}$ calcula: (a) el terme quart, (b) el terme independent, i (c) el coeficient del terme de lloc 14.
\end{exercici}

\begin{solucio}
Recorda que el terme $k$-èsim del desenvolupament $(A+B)^{n}$ ve
donat per\[
t_{k}=\binom{n}{k-1}\cdot A^{n-(k-1)}\cdot B^{k-1}
\]

a)
\begin{eqnarray*}
t_{4} &=&\binom{15}{4-1}\cdot \left( 2x\right) ^{15-(4-1)}\cdot \left( \frac{1}{2}\right) ^{4-1} \\
&=&\binom{15}{3}\cdot 2^{12}\cdot x^{12}\cdot \frac{1}{2^{3}} \\
&=&455\cdot 2^{9}\cdot x^{12} \\
&=&232960\cdot x^{12}
\end{eqnarray*}perquè,\[
\binom{15}{3}=\frac{15!}{3!\cdot (15-3)!}=455
\]

b) El terme independent és aquell que no conté $x$. Per tant,
com el terme general és donat per l'expressió\[
t_{k}=\binom{15}{k-1}\cdot \left( 2x\right) ^{15-(k-1)}\cdot \left( \frac{1}{2}\right) ^{k-1}
\]cal complir-se\begin{eqnarray*}
15-(k-1) &=&0 \\
k &=&16
\end{eqnarray*}Com era d'esperar, en aquest cas el terme independent és l'últim
terme. Aleshores,\begin{eqnarray*}
t_{16} &=&\binom{15}{15}\cdot \left( 2x\right) ^{0}\cdot \left( \frac{1}{2}\right) ^{15} \\
&=&\frac{1}{2^{15}}
\end{eqnarray*}

c) El terme de lloc 14 vindrà donat per\begin{eqnarray*}
t_{14} &=&\binom{15}{14-1}\cdot \left( 2x\right) ^{15-(14-1)}\cdot \left(
\frac{1}{2}\right) ^{14-1} \\
&=&\binom{15}{13}\cdot \left( 2x\right) ^{2}\cdot \left( \frac{1}{2}\right)
^{13} \\
&=&\binom{15}{13}\cdot \frac{1}{2^{11}}\cdot x^{2}
\end{eqnarray*}i, per tant, el seu coeficient és\[
\binom{15}{13}\cdot \frac{1}{2^{11}}=\frac{105}{2048}
\]
\end{solucio}

\begin{exercici}
De la potència del binomi $\left( 3x-\frac{1}{3x}\right) ^{6}$ escriu el terme general del seu desenvolupament per la fórmula de Newton. Troba el terme tercer i el central.
\end{exercici}

\begin{solucio}
El terme $k$-èsim del desenvolupament $(A+B)^{n}$ ve donat per\[
t_{k}=\binom{n}{k-1}\cdot A^{n-(k-1)}\cdot B^{k-1}
\]Aleshores, en el nostre cas, és\begin{eqnarray*}
t_{k} &=&\binom{6}{k-1}\cdot \left( 3x\right) ^{6-(k-1)}\cdot \left( -\frac{1}{3x}\right) ^{k-1} \\
&=&(-1)^{k-1}\cdot \binom{6}{k-1}\cdot \left( 3x\right) ^{7-k}\cdot \frac{1}{\left( 3x\right) ^{k-1}} \\
&=&(-1)^{k-1}\cdot \binom{6}{k-1}\cdot \left( 3x\right) ^{8-2k}
\end{eqnarray*}

Aleshores, el terme tercer és\begin{eqnarray*}
t_{3} &=&(-1)^{3-1}\cdot \binom{6}{3-1}\cdot \left( 3x\right) ^{8-6} \\
&=&\binom{6}{2}\cdot \left( 3x\right) ^{2} \\
&=&135x^{2}
\end{eqnarray*}

Com que el desenvolupament conté set termes, el terme central
és el quart terme. Per tant,
\begin{eqnarray*}
t_{4} &=&(-1)^{4-1}\cdot \binom{6}{4-1}\cdot \left( 3x\right) ^{8-8} \\
&=&-\binom{6}{3} \\
&=&-20
\end{eqnarray*}
\end{solucio}

\begin{exercici}
Del desenvolupament de $(x-1)^{200}$ calcula: (a) el terme de lloc 198,
(b) l'últim terme, i (c) el terme en què apareix $x^{198}$.
\end{exercici}

\begin{solucio}
El terme $k$-èsim del desenvolupament $(A+B)^{n}$ ve donat per\[
t_{k}=\binom{n}{k-1}\cdot A^{n-(k-1)}\cdot B^{k-1}
\]

a) Aleshores, el terme de lloc 198 és\begin{eqnarray*}
t_{198} &=&\binom{200}{197}\cdot x^{200-197}\cdot \left( -1\right) ^{197} \\
&=&-\binom{200}{197}\cdot x^{3}
\end{eqnarray*}

b) L'últim terme és el de lloc $201$ i, per tant, tenim\begin{eqnarray*}
t_{201} &=&\binom{200}{200}\cdot x^{200-200}\cdot \left( -1\right) ^{200} \\
&=&\binom{200}{200}=1
\end{eqnarray*}

c) Per a calcular el terme en què apareix $x^{198}$, primer,
calculem el terme general\[
t_{k}=\binom{200}{k-1}\cdot x^{200-(k-1)}\cdot \left( -1\right) ^{k-1}
\]Aleshores, s'ha de complir que\begin{eqnarray*}
200-(k-1) &=&198 \\
k &=&3
\end{eqnarray*}Aleshores, el terme demanat és\begin{eqnarray*}
t_{3} &=&\binom{200}{2}\cdot x^{200-2}\cdot \left( -1\right) ^{2} \\
&=&\binom{200}{2}\cdot x^{198}
\end{eqnarray*}
\end{solucio}

\begin{exercici}
Troba el terme de lloc 16 en el desenvolupament de $\left( \sqrt{3}-\sqrt[5]{3}\right) ^{19}$, expressant-lo de la manera més senzill possible.
\end{exercici}

\begin{solucio}
El terme $k$-èsim del desenvolupament $(A+B)^{n}$ ve donat per\[
t_{k}=\binom{n}{k-1}\cdot A^{n-(k-1)}\cdot B^{k-1}
\]Per tant, en el nostre cas, obtenim\begin{eqnarray*}
t_{16} &=&\binom{19}{15}\cdot \left( \sqrt{3}\right) ^{19-15}\cdot \left( -\sqrt[5]{3}\right) ^{15} \\
&=&-\binom{19}{15}\cdot \left( \sqrt{3}\right) ^{4}\cdot \left( \sqrt[5]{3}\right) ^{15} \\
&=&-\binom{19}{15}\cdot 3^{2}\cdot 3^{3} \\
&=&-\binom{19}{15}\cdot 3^{5}
\end{eqnarray*}
\end{solucio}

\section{Divisió de polinomis. Algorisme d’Euclides}

\begin{exercici}
Calcula el quocient i el residu de les divisions següents\[
\begin{array}{l}
a)~(3x^{2}+2x-7):(x^{2}+3) \\
b)~(x^{3}+5x-6):(2x+3) \\
c)~(4x^{4}-1):(3x^{2}+5x+3) \\
d)~\left( \frac{3}{4}x^{4}-\frac{1}{3}x^{3}-7\right):(x^{2}+x-1) \\
e)~(x^{5}-1):(x+1)\end{array}\]
\end{exercici}

\begin{solucio}
a) Observa que\[
\begin{array}{ccccc|ccc}
3x^{2} & + & 2x & - & 7 & x^{2} & + & 3 \\ \cline{6-8}
-3x^{2} & & & - & 9 & 3 & & \\ \cline{1-5}
& \diagup & 2x & - & 16 & & &
\end{array}\]i, per tant, el quocient és $3$, i el residu, $2x-16$.

b) De la mateixa manera,\[
\begin{array}{ccccccc|ccccc}
x^{3} & + & 0x^{2} & + & 5x & - & 6 & 2x & + & 3 & & \\ \cline{8-12}
-x^{3} & - & \frac{3}{2}x^{2} & & & & & \frac{1}{2}x^{2} & - & \frac{3}{4}x & + & \frac{29}{8} \\ \cline{1-5}
\diagup & - & \frac{3}{2}x^{2} & + & 5x & & & & & & & \\
& & \frac{3}{2}x^{2} & + & \frac{9}{4}x & & & & & & & \\ \cline{3-7}
& & \diagup & & \frac{29}{4}x & - & 6 & & & & & \\
& & & & -\frac{29}{4}x & - & \frac{87}{8} & & & & & \\ \cline{5-7}
& & & & \diagup & - & \frac{135}{8} & & & & &
\end{array}\]i, per tant, el quocient és\[
\frac{1}{2}x^{2}-\frac{3}{4}x+\frac{29}{8}
\]i el residu, és $-135/8$.

c) Tenim,\[
\begin{array}{ccccccccc|ccccc}
4x^{4} & + & 0x^{3} & + & 0x^{2} & + & 0x & - & 1 & 3x^{2} & + & 5x & + & 3
\\ \cline{10-14}
-4x^{4} & - & \frac{20}{3}x^{3} & - & 4x^{2} & & & & & \frac{4}{3}x^{2}
& - & \frac{20}{9}x & + & \frac{64}{27} \\ \cline{1-5}
\diagup & - & \frac{20}{3}x^{3} & - & 4x^{2} & & & & & & & & & \\
& & \frac{20}{3}x^{3} & + & \frac{100}{9}x^{2} & + & \frac{20}{3}x & & &
& & & & \\ \cline{3-7}
& & \diagup & & \frac{64}{9}x^{2} & + & \frac{20}{3}x & - & 1 & & & &
& \\
& & & - & \frac{64}{9}x^{2} & - & \frac{320}{27}x & - & \frac{64}{9} & &
& & & \\ \cline{5-9}
& & & & \diagup & - & \frac{140}{27}x & - & \frac{73}{9} & & & & &
\end{array}\]i, per tant, el quocient és\[
\frac{4}{3}x^{2}-\frac{20}{9}x+\frac{64}{27}
\]i el residu és
\[
-\frac{140}{27}x-\frac{73}{9}
\]

d) Tenim\[
\begin{array}{ccccccccc|ccccc}
\frac{3}{4}x^{4} & - & \frac{1}{3}x^{3} & + & 0x^{2} & + & 0x & - & 7 & x^{2}
& + & x & - & 1 \\ \cline{10-14}
-\frac{3}{4}x^{4} & - & \frac{3}{4}x^{3} & + & \frac{3}{4}x^{2} & & & &
& \frac{3}{4}x^{2} & - & \frac{13}{12}x & + & \frac{11}{6} \\ \cline{1-5}
\diagup & - & \frac{13}{12}x^{3} & + & \frac{3}{4}x^{2} & & & & & & &
& & \\
& & \frac{13}{12}x^{3} & + & \frac{13}{12}x^{2} & - & \frac{13}{12}x & &
& & & & & \\ \cline{3-7}
& & \diagup & & \frac{11}{6}x^{2} & - & \frac{13}{12}x & - & 7 & & & &
& \\
& & & - & \frac{11}{6}x^{2} & - & \frac{11}{6}x & + & \frac{11}{6} & & &
& & \\ \cline{5-9}
& & & & \diagup & - & \frac{35}{12}x & - & \frac{31}{6} & & & & &
\end{array}\]i, per tant, el quocient és\[
\frac{3}{4}x^{2}-\frac{13}{12}x+\frac{11}{6}
\]i el residu,\[
-\frac{35}{12}x-\frac{31}{6}
\]

e) Tenim\[
\begin{array}{ccccccccccc|ccccccccc}
x^{5} & + & 0x^{4} & + & 0x^{3} & + & 0x^{2} & + & 0x & - & 1 & x & + & 1 &
& & & & & \\ \cline{12-20}
-x^{5} & - & x^{4} & & & & & & & & & x^{4} & - & x^{3} & + & x^{2} &
- & x & + & 1 \\ \cline{1-3}
\diagup & - & x^{4} & & & & & & & & & & & & & & & & & \\
& & x^{4} & + & x^{3} & & & & & & & & & & & & & & & \\
\cline{3-5}
& & \diagup & & x^{3} & & & & & & & & & & & & & & & \\
& & & - & x^{3} & - & x^{2} & & & & & & & & & & & & & \\
\cline{5-7}
& & & & \diagup & - & x^{2} & & & & & & & & & & & & & \\
& & & & & & x^{2} & + & x & & & & & & & & & & & \\
\cline{7-9}
& & & & & & \diagup & & x & - & 1 & & & & & & & & & \\
& & & & & & & - & x & - & 1 & & & & & & & & & \\ \cline{9-11}
& & & & & & & & \diagup & - & 2 & & & & & & & & &
\end{array}\]i, per tant, el quocient és\[
x^{4}-x^{3}+x^{2}-x+1
\]i el residu, $-2$.
\end{solucio}

\begin{exercici}
Troba un polinomi de tercer grau que dividit per $x+1$ dé de quocient
$x^{2}+3$ i residu $-2$.
\end{exercici}

\begin{solucio}
Si $p(x)$ és el polinomi que dividit per $x+1$ dona com a quocient $x^{2}+3$
i com a residu $-2$, aleshores s'ha de complir la següent condició:\begin{eqnarray*}
p(x) &=&(x+1)\cdot (x^{2}+3)-2 \\
&=&x^{3}+3x+x^{2}+3-2 \\
&=&x^{3}+x^{2}+3x+1
\end{eqnarray*}
\end{solucio}

\begin{exercici}
Esbrinar si les divisions següents són o no exactes:\[
\begin{array}{l}
a)~(x^{2}+2x+1):(x+1) \\
b)~(x^{3}-3x^{2}+x-3):(x^{2}+1) \\
c)~\left( x^{3}+2x^{2}+\frac{15}{4}x+\frac{9}{4}\right):\left( x^{2}+\frac{1}{2}x+3\right)\end{array}\]
\end{exercici}

\begin{solucio}
Un mètode per a esbrinar si la divisió és exacta o no és efectuarla
i veure si el residu és zero o no. Tanmateix, aquí no procedirem així i, en el seu lloc, farem servir altres alternatives més senzilles.

a) Es evident que la divisió de $x^{2}+2x+1$ per $x+1$ és exacta,
perquè, $x^{2}+2x+1$ és un múltiple de $x+1$ al complir-se\[
x^{2}+2x+1=(x+1)^{2}=(x+1)\cdot (x+1)
\]

b) Si la divisió de $x^{3}-3x^{2}+x-3$ per $x^{2}+1$ fos exacta,
existiria un polinomi de primer grau (quocient) $c(x)=ax+b$ tal que\begin{eqnarray*}
x^{3}-3x^{2}+x-3 &=&(x^{2}+1)(ax+b) \\
&=&ax^{3}+bx^{2}+ax+b
\end{eqnarray*}i, per tant, caldria complir-se\[
a=1\text{ \ \ i \ \ }b=-3
\]Com que existeixen aquests coeficients, podem assegurar que la divisió és
exacta i, a més, es compleix\[
x^{3}-3x^{2}+x-3=(x^{2}+1)\cdot (x-3)
\]

c) Anàlogament, si la divisió de\[
x^{3}+2x^{2}+\frac{15}{4}x+\frac{9}{4}
\]
per\[
x^{2}+\frac{1}{2}x+3
\]fos exacta, existiria un polinomi de primer grau (quocient) $c(x)=ax+b$ tal que\begin{eqnarray*}
x^{3}+2x^{2}+\frac{15}{4}x+\frac{9}{4} &=&\left( x^{2}+\frac{1}{2}x+3\right)
\cdot \left( ax+b\right) \\
&=&ax^{3}+\frac{1}{2}x^{2}+3ax+bx^{2}+\frac{b}{2}x+3b \\
&=&ax^{3}+\left( \frac{1}{2}+b\right) x^{2}+\left( 3a+\frac{b}{2}\right) x+3b
\end{eqnarray*}i, per tant, caldria complir-se simultàniament les següents
equacions\[
\begin{array}{c}
a=1 \\
\frac{1}{2}+b=2 \\
3a+\frac{b}{2}=\frac{15}{4} \\
3b=\frac{9}{4}\end{array}\]Ara bé, la segona i la quarta són incompatibles, perquè, de la segona es
dedueix $b=3/2$ i, de la quarta, $b=3/4$, resultando dos valors diferents
per a $b$. Per consegüent, la divisió en aquest cas no és exacta.
\end{solucio}

\begin{exercici}
Calcula $a$ perquè $x^{3}-3x^{2}+x-a$ sigui divisible per $x^{2}+1$.
\end{exercici}

\begin{solucio}
Si $x^{3}-3x^{2}+x-a$ és divisible per $x^{2}+1$ significa que la divisió de $x^{3}-3x^{2}+x-a$ per $x^{2}+1$ és exacta i, en conseqüència, existe
un polinomi de primer grau (quocient) $c(x)=mx+n$ tal que\begin{eqnarray*}
x^{3}-3x^{2}+x-a &=&(x^{2}+1)\cdot (mx+n) \\
&=&mx^{3}+nx^{2}+mx+n
\end{eqnarray*}Aleshores, s'han de satisfer conjuntamente les següents condicions\[
\begin{array}{c}
m=1 \\
n=-3 \\
n=-a\end{array}\]i, per tant, $a=3$.
\end{solucio}

\begin{exercici}
Donat el polinomi $3x^{5}+7x^{2}-5$, troba el seu polinomi mònic.
\end{exercici}

\begin{solucio}
Recorda que un polinomi s’anomena mònic si el coeficient del monomi
de grau més gran és igual a la unitat. És clar que\[
3x^{5}+7x^{2}-5=3\left( x^{5}+\frac{7}{3}x^{2}-\frac{5}{3}\right)
\]i, per tant,\[
x^{5}+\frac{7}{3}x^{2}-\frac{5}{3}
\]és el polinomi mònic de $3x^{5}+7x^{2}-5$.
\end{solucio}

\begin{exercici}
Utilitzant l'algorisme d'Euclides, calcula m.c.d. i m.c.m. dels
polinomis $p(x)$ i $q(x)$, en els casos següents:\[
\begin{array}{l}
a)~p(x)=x^{4}-5x^{2}+6\text{ i }q(x)=x^{3}+x^{2}-3x-3 \\
b)~p(x)=x^{6}+x^{4}-x^{2}-1\text{ i }q(x)=x^{5}-x^{4}-x+1 \\
c)~p(x)=x^{3}+\sqrt{2}x^{2}+3x+3\sqrt{2}\text{ i }q(x)=x^{2}+2\end{array}\]
\end{exercici}

\begin{solucio}
Recorda que l'algorisme d'Euclides consisteix a fer divisions successives, de manera que\begin{eqnarray*}
\text{m.c.d.}\left( p(x),q(x)\right) &=&\text{m.c.d.}\left(
q(x),r_{1}(x)\right) \\
&=&\text{m.c.d.}\left( r_{1}(x),r_{2}(x)\right) \\
&&\cdots \\
&=&\text{m.c.d.}\left( r_{n-1}(x),r_{n}(x)\right) \\
&=&r_{n-1}^{\ast }(x)
\end{eqnarray*}essent $r_{n}(x)$ un nombre real i $r_{n-1}^{\ast }(x)$ el polinomi mònic de $r_{n-1}(x)$. Un cop calculat el m.c.d., calcularem m.c.m.
mitjançant\[
\text{m.c.m.}\left( p(x),q(x)\right) =\frac{1}{a_{n}\cdot b_{m}}\frac{p(x)\cdot q(x)}{\text{m.c.d.}\left( p(x),q(x)\right) }
\]

a)
\[
\begin{array}{ccccc|cccc|ccc}
& & & & & x & -1 & & & x & +1 & \\ \cline{6-12}
x^{4} & 0 & -5x^{2} & 0 & +6 & x^{3} & +x^{2} & -3x & -3 & x^{2} & -3 & \\
-x^{4} & -x^{3} & +3x^{2} & +3x & & -x^{3} & & +3x & & & & \\
\cline{1-4}\cline{6-8}
\diagup & -x^{3} & -2x^{2} & +3x & +6 & \diagup & x^{2} & 0 & -3 & & & \\
& x^{3} & x^{2} & -3x & -3 & & -x^{2} & & +3 & & & \\
\cline{2-5}\cline{7-9}
& \diagup & -x^{2} & 0 & 3 & & \diagup & & \diagup & & & \\
& & x^{2} & & -3 & & & & & & &
\end{array}\]Per tant, tenim\begin{eqnarray*}
\text{m.c.d.}\left( x^{4}-5x^{2}+6,x^{3}+x^{2}-3x-3\right) &=&\text{m.c.d.}\left( x^{3}+x^{2}-3x-3,x^{2}-3\right) \\
&=&\text{m.c.d.}\left( x^{2}-3,0\right) \\
&=&x^{2}-3
\end{eqnarray*}i\begin{eqnarray*}
\text{m.c.m.}\left( x^{4}-5x^{2}+6,x^{3}+x^{2}-3x-3\right) &=&\frac{\left(
x^{4}-5x^{2}+6\right) \cdot \left( x^{3}+x^{2}-3x-3\right) }{x^{2}-3} \\
&=&x^{5}-5x^{3}+x^{4}-5x^{2}+6x+6
\end{eqnarray*}

b)
\[
\begin{array}{ccccccc|cccccc|cc}
& & & & & & & x & +1 & & & & & x & -1 \\ \cline{8-15}
x^{6} & 0 & +x^{4} & 0 & -x^{2} & 0 & -1 & x^{5} & -x^{4} & 0 & 0 & -x & +1
& x^{4} & -1 \\
-x^{6} & +x^{5} & & & +x^{2} & -x & & -x^{5} & & & & +x & & & \\
\cline{1-6}\cline{8-12}
\diagup & x^{5} & +x^{4} & 0 & & -x & -1 & \diagup & -x^{4} & 0 & 0 & 0 & 1
& & \\
& -x^{5} & +x^{4} & & & +x & -1 & & x^{4} & & & & -1 & & \\
\cline{2-7}\cline{9-13}
& \diagup & 2x^{4} & 0 & 0 & 0 & -2 & & \diagup & & & & 0 & & \\
& & x^{4} & & & & -1 & & & & & & & &
\end{array}\]Aleshores,\begin{eqnarray*}
\text{m.c.d.}\left( x^{6}+x^{4}-x^{2}-1,x^{5}-x^{4}-x+1\right) &=&\text{m.c.d.}\left( x^{3}+x^{2}-3x-3,x^{5}-x^{4}-x+1\right) \\
&=&\text{m.c.d.}\left( x^{5}-x^{4}-x+1,x^{4}-1\right) \\
&=&\text{m.c.d.}\left( x^{4}-1,0\right) \\
&=&x^{4}-1
\end{eqnarray*}i\begin{eqnarray*}
\text{m.c.m.}\left( x^{6}+x^{4}-x^{2}-1,x^{5}-x^{4}-x+1\right) &=&\frac{\left( x^{6}+x^{4}-x^{2}-1\right) \cdot \left( x^{5}-x^{4}-x+1\right) }{x^{4}-1} \\
&=&x^{7}+x^{5}-x^{6}-x^{4}-x^{3}-x+x^{2}+1
\end{eqnarray*}

c)\[
\begin{array}{cccc|ccc|cc|cc}
& & & & x & +\sqrt{2} & & x & -\sqrt{2} & x & +\sqrt{2}
\\\cline{5-11}
x^{3} & +\sqrt{2}x^{2} & +3x & +3\sqrt{2} & x^{2} &
0 & +2 & x & +\sqrt{2} &
1 & \\
-x^{3} & & -2x & & -x^{2} & -\sqrt{2}x & & -x & -\sqrt{2} & & \\
\cline{1-3}\cline{5-6}\cline{8-9}
\diagup & \sqrt{2}x^{2} & +x & +3\sqrt{2} & \diagup & -\sqrt{2}x & +2 &
\diagup & 0 & & \\
& -\sqrt{2}x^{2} & & -2\sqrt{2} & & \sqrt{2}x & +2 & & & & \\
\cline{2-4}\cline{6-7}
& \diagup & x & +\sqrt{2} & & \diagup & 4 & & & & \\
& & & & & & 1 & & & &
\end{array}\]Aleshores, tenim\begin{eqnarray*}
\text{m.c.d.}\left( x^{3}+\sqrt{2}x^{2}+3x+3\sqrt{2},x^{2}+2\right) &=&\text{m.c.d.}\left( x^{2}+2,x+\sqrt{2}\right) \\
&=&\text{m.c.d.}\left( x+\sqrt{2},1\right) \\
&=&1
\end{eqnarray*}Observa que els dos polinomis són primers entre si i, per tant,\begin{eqnarray*}
\text{m.c.m.}\left( x^{3}+\sqrt{2}x^{2}+3x+3\sqrt{2},x^{2}+2\right)
&=&\left( x^{3}+\sqrt{2}x^{2}+3x+3\sqrt{2}\right) \cdot \left( x^{2}+2\right)
\\
&=&x^{5}+5x^{3}+\sqrt{2}x^{4}+5\sqrt{2}x^{2}+6x+6\sqrt{2}
\end{eqnarray*}
\end{solucio}

\begin{exercici}
Calcula $a$ i $b$ perquè el polinomi m.c.d. de $p(x)=x^{2}+ax+b$ i $q(x)=x^{2}-9$ sigui $x+3$. Hi ha una única solució?
\end{exercici}

\begin{solucio}
Segons el enunciat, tenim que
\[
\text{m.c.d.}\left( x^{2}+ax+b,x^{2}-9\right) =x+3
\]Per tant, $x^{2}+ax+b$ i $x^{2}-9$ han de ser divisibles per $x+3$. Es
és clar que\[
x^{2}-9=(x+3)\cdot (x-3)
\]i, per tant, cal verificar-se\begin{eqnarray*}
x^{2}+ax+b &=&(x+3)\cdot (x+k) \\
&=&x^{2}+(k+3)x+3k
\end{eqnarray*}i, d'aquí, es dedueix que s'han de complir les relacions següents\begin{eqnarray*}
k+3 &=&a \\
3k &=&b
\end{eqnarray*}és a dir,\[
3(a-3)=b
\]Per consegüent, hi ha més d’una solució. Per exemple, si fem $a=1$, aleshores $b=-6$, i es compleix\[
\text{m.c.d.}\left( x^{2}+x-6,x^{2}-9\right) =x+3
\]
\end{solucio}

\begin{exercici}
 És divisible $p(x)=2x^{3}+3x^{2}+2x+1$ per $q(x)=x^{3}+1$?
Calcula el m.c.d. i el m.c.m. de tots dos polinomis.
\end{exercici}

\begin{solucio}
Com que és $p(x)$ i $q(x)$ del mateix grau, és evident que $p(x)$ no és
divisible per $q(x)$. Per l'algorisme d'Euclides, tenim\[
\begin{array}{cccc|cccc|ccc|cc}
& & & & 2 & & & & x & -\frac{2}{3} & & x & -\frac{1}{3} \\
\cline{5-13}
2x^{3} & +3x^{2} & +2x & +1 & x^{3} & 0 & 0 & 1 & x^{2} & +\frac{2}{3}x & -\frac{1}{3} & x & +1 \\
-2x^{3} & & & -2 & -x^{3} & -\frac{2}{3}x^{2} & +\frac{1}{3}x & & -x^{2}
& -x & & & \\ \cline{1-4}\cline{5-7}\cline{9-10}
\diagup & 3x^{2} & +2x & -1 & \diagup & -\frac{2}{3}x^{2} & +\frac{1}{3}x & 1
& \diagup & -\frac{1}{3}x & -\frac{1}{3} & & \\
& x^{2} & +\frac{2}{3}x & -\frac{1}{3} & & \frac{2}{3}x^{2} & +\frac{4}{9}x
& -\frac{2}{9} & & \frac{1}{3}x & +\frac{1}{3} & & \\
\cline{6-8}\cline{10-11}
& & & & & \diagup & \frac{7}{9}x & +\frac{7}{9} & & \diagup & 0 & &
\\
& & & & & & x & +1 & & & & &
\end{array}\]Per tant, tenim\begin{eqnarray*}
\text{m.c.d.}\left( 2x^{3}+3x^{2}+2x+1,x^{3}+1\right) &=&\text{m.c.d.}\left(
x^{3}+1,x^{2}+\frac{2}{3}x-\frac{1}{3}\right) \\
&=&\text{m.c.d.}\left( x^{2}+\frac{2}{3}x-\frac{1}{3},x+1\right) \\
&=&\text{m.c.d.}\left( x+1,0\right) \\
&=&x+1
\end{eqnarray*}D'altra banda, tenim\begin{eqnarray*}
\text{m.c.m.}\left( 2x^{3}+3x^{2}+2x+1,x^{3}+1\right) &=&\frac{1}{2}\frac{\left( 2x^{3}+3x^{2}+2x+1\right) \left( x^{3}+1\right) }{x+1} \\
&=&x^{5}+\frac{1}{2}x^{4}+\frac{1}{2}x^{3}+x^{2}+\frac{1}{2}x+\frac{1}{2}
\end{eqnarray*}
\end{solucio}

\section{Regla de Ruffini i teorema del residu}

\begin{exercici}
Utiliza la regla de Ruffini per a determinar el quocient i el residu de les
següents divisions:\[
\begin{array}{l}
a)~(x^{5}-x^{3}-x-1):(x-3) \\
b)~(x^{3}-4):\left( x-\frac{1}{2}\right) \\
c)~(x^{4}-3):(x-\sqrt{3}) \\
d)~(x^{3}-2x^{2}+5):\left( x-\frac{1}{3}\right)\end{array}\]
\end{exercici}

\begin{solucio}
a) Mitjançant la regla de Ruffini, tenim\[
\begin{array}{r|rrrrrr}
& 1 & 0 & -1 & 0 & -1 & -1 \\
3 & & 3 & 9 & 24 & 72 & 213 \\ \hline
& 1 & 3 & 8 & 24 & 71 & 212\end{array}\]i, per tant, el quocient és $c(x)=x^{4}+3x^{3}+8x^{2}+24x+71$, i el residu, $r(x)=212$.

b) Mitjançant la regla de Ruffini, tenim\[
\begin{array}{r|rrrr}
& 1 & 0 & 0 & -4 \\
\frac{1}{2} & & \frac{1}{2} & \frac{1}{4} & \frac{1}{8} \\ \hline
& 1 & \frac{1}{2} & \frac{1}{4} & -\frac{31}{8}\end{array}\]i, per tant, el quocient és $c(x)=x^{2}+\frac{1}{2}x+\frac{1}{4}$, i el
residu, $r(x)=-\frac{31}{8}$.

c) Mitjançant la regla de Ruffini, tenim\[
\begin{array}{r|rrrrr}
& 1 & 0 & 0 & 0 & -3 \\
\sqrt{3} & & \sqrt{3} & 3 & 3\sqrt{3} & 9 \\ \hline
& 1 & \sqrt{3} & 3 & 3\sqrt{3} & 6\end{array}\]i, per tant, el quocient és $c(x)=x^{3}+\sqrt{3}x^{2}+3x+3\sqrt{3}$, i el
residu, $r(x)=6$.

d) Mitjançant la regla de Ruffini, tenim\[
\begin{array}{r|rrrr}
& 1 & -2 & 0 & 5 \\
\frac{1}{3} & & \frac{1}{3} & -\frac{5}{9} & -\frac{5}{27} \\ \hline
& 1 & -\frac{5}{3} & -\frac{5}{9} & \frac{130}{27}\end{array}\]i, per tant, el quocient és $c(x)=x^{2}-\frac{5}{3}x-\frac{5}{9}$, i el
residu, $r(x)=\frac{130}{27}$.
\end{solucio}

\begin{exercici}
Utilizal teorema del residu per a determinar el residu de cada una de les
següents divisions:\[
\begin{array}{l}
a)~(x^{2}+x+1):\left( x-\frac{1}{3}\right) \\
b)~(x^{3}-2x+\sqrt{2}):(x-\sqrt{2}) \\
c)~\left( \frac{3}{4}x^{5}+8x-\frac{5}{7}\right):\left( x+\frac{2}{3}\right)
\\
d)~(4x^{4}+5x^{3}-7x^{2}):(x+2)\end{array}\]
\end{exercici}

\begin{solucio}
Recorda que, segons el teorema del residu, el valor numèric del
polinomi $p(x)$ per a $x=a$ coincideix amb el residu de dividir $p(x)$ per $x-a$.

a) Si fem $p(x)=x^{2}+x+1$, aleshores el residu de la divisió per $x-\frac{1}{3}$ és\begin{eqnarray*}
p\left( \frac{1}{3}\right) &=&\left( \frac{1}{3}\right) ^{2}+\frac{1}{3}+1 \\
&=&\frac{13}{9}
\end{eqnarray*}

b) Si fem $p(x)=x^{3}-2x+\sqrt{2}$, aleshores el residu de la divisió
per $x-\sqrt{2}$ és\begin{eqnarray*}
p\left( \sqrt{2}\right) &=&\left( \sqrt{2}\right) ^{3}-2\sqrt{2}+\sqrt{2} \\
&=&2\sqrt{2}-2\sqrt{2}+\sqrt{2} \\
&=&\sqrt{2}
\end{eqnarray*}

c) Si fem $p(x)=\frac{3}{4}x^{5}+8x-\frac{5}{7}$, aleshores el residu de
la divisió per $x+\frac{2}{3}$ és\begin{eqnarray*}
p\left( -\frac{2}{3}\right) &=&\frac{3}{4}\left( -\frac{2}{3}\right)
^{5}+8\left( -\frac{2}{3}\right) -\frac{5}{7} \\
&=&-\frac{8}{81}-\frac{16}{3}-\frac{5}{7} \\
&=&-\frac{3485}{567}
\end{eqnarray*}

d) Si fem $p(x)=4x^{4}+5x^{3}-7x^{2}$, aleshores el residu de la divisió per $x+2$ és\begin{eqnarray*}
p(-2) &=&4\left( -2\right) ^{4}+5\left( -2\right) ^{3}-7\left( -2\right) ^{2}
\\
&=&64-40-28 \\
&=&-4
\end{eqnarray*}
\end{solucio}

\begin{exercici}
Sense efectuar la divisió, averigua si els polinomis que siguen són
divisibles pels polinomis indicats\[
\begin{array}{l}
a)~x^{4}-x^{2}+6\text{ i }x+\sqrt{3} \\
b)~x^{2}-x-1\text{ i }x-\frac{1+\sqrt{5}}{2} \\
c)~x^{6}-x^{4}-x^{2}-1\text{ i }x-1\end{array}\]
\end{exercici}

\begin{solucio}
a) El polinomi $p(x)=x^{4}-x^{2}+6$ serà divisible per $x+\sqrt{3}$, si
el residu de la divisió és zero. Aplicando el teorema del residu, tenim\begin{eqnarray*}
p\left( -\sqrt{3}\right) &=&\left( -\sqrt{3}\right) ^{4}-\left( -\sqrt{3}\right) ^{2}+6 \\
&=&9-3+6 \\
&=&12
\end{eqnarray*}i, com a conseqüència, $p(x)$ no és divisible per $x+\sqrt{3}$.

b) El polinomi $p(x)=x^{2}-x-1$ serà divisible per $x-\frac{1+\sqrt{5}}{2}$, si el residu de la divisió és zero. Aplicando el teorema del residu,
tenim\begin{eqnarray*}
p\left( \frac{1+\sqrt{5}}{2}\right) &=&\left( \frac{1+\sqrt{5}}{2}\right)
^{2}-\frac{1+\sqrt{5}}{2}-1 \\
&=&\frac{1+2\sqrt{5}+5}{4}-\frac{1+\sqrt{5}}{2}-1 \\
&=&\frac{6+2\sqrt{5}-2-2\sqrt{5}-4}{4} \\
&=&0
\end{eqnarray*}i, per tant, el polinomi és divisible per $x-\frac{1+\sqrt{5}}{2}$.

c) El polinomi $p(x)=x^{6}-x^{4}-x^{2}-1$ serà divisible per $x-1$, si
el residu de la divisió és zero. Aplicando el teorema del residu, tenim\begin{eqnarray*}
p(1) &=&1^{6}-1^{4}-1^{2}-1 \\
&=&-2
\end{eqnarray*}i, com a conseqüència, $p(x)$ no és divisible per $x-1$.
\end{solucio}

\begin{exercici}
Escriu un polinomi de tercer grau que sigui divisible pels polinomis $x-3+\frac{\sqrt{2}}{2}$ i $x-3-\frac{\sqrt{2}}{2}$.
\end{exercici}

\begin{solucio}
Si el polinomi que busquem $p(x)$ és divisible per $x-3+\frac{\sqrt{2}}{2}$, aleshores\[
p(x)=\left( x-3+\frac{\sqrt{2}}{2}\right) \cdot c_{1}(x)
\]on $c_{1}(x)$ és un polinomi de segon grau. Si, a més, $p(x)$ és
divisible per $x-3-\frac{\sqrt{2}}{2}$, aleshores pel teorema del residu,
obtenim\[
c_{1}(x)=\left( x-3-\frac{\sqrt{2}}{2}\right) \cdot c_{2}(x)
\]on $c_{2}(x)$ és un polinomi de primer grau. Per tant, es compleix\[
p(x)=\left( x-3+\frac{\sqrt{2}}{2}\right) \cdot \left( x-3-\frac{\sqrt{2}}{2}\right) \cdot c_{2}(x)
\]Si ara prenem $c_{2}(x)=x$, obtenim\[
p(x)=x^{3}-6x^{2}+\frac{17}{2}x
\]que compleix les condicions de l'enunciat.
\end{solucio}

\begin{exercici}
Calcula els valors de $a$ i $b$ perquè el polinomi $p(x)=x^{4}-5x^{3}+4x^{2}+ax-b$ sigui divisible pels binomis $x-3$ i $x+1$.
\end{exercici}

\begin{solucio}
Si $p(x)$ és divisible per $x-3$, aleshores pel teorema del residu,
obtenim\[
p(3)=3^{4}-5\cdot 3^{3}+4\cdot 3^{2}+a\cdot 3-b=0
\]o sigui,\begin{eqnarray*}
81-135+36+3a-b &=&0 \\
3a-b &=&18
\end{eqnarray*}Si, a més, $p(x)$ és divisible per $x+1$, aleshores pel teorema del
residu, obtenim\[
p(-1)=\left( -1\right) ^{4}-5\left( -1\right) ^{3}+4\left( -1\right)
^{2}+a\left( -1\right) -b=0
\]és a dir,\begin{eqnarray*}
1+5+4-a-b &=&0 \\
-a-b &=&-10 \\
a+b &=&10
\end{eqnarray*}En resum, perquè es compleixin les condicions de l'enunciat s'han de
satisfer conjuntamente les equacions\begin{eqnarray*}
3a-b &=&18 \\
a+b &=&10
\end{eqnarray*}Resolent aquest sistema, obtenim $a=7$ i $b=3$.
\end{solucio}

\begin{exercici}
Troba el valor de $a$ perquè el polinomi $p(x)=x^{3}+ax+3$ doni de
residu $-4$ en dividir-lo per $x-1$.
\end{exercici}

\begin{solucio}
Recorda que, segons el teorema del residu, el valor numèric del
polinomi $p(x)$ per a $x=a$ coincideix amb el residu de dividir $p(x)$ per $x-a$. Per tant, s'ha de verificar que\[
p(1)=-4
\]és a dir,\begin{eqnarray*}
1+a+3 &=&-4 \\
a &=&-8
\end{eqnarray*}
\end{solucio}

\begin{exercici}
Se sap que un polinomi $p(x)$ és divisible per $x+1$ i, a més, el
residu de la divisió de $p(x)$ per $x+2$ és $4$. Calcula el residu que es
obté en dividir-lo per $(x+1)(x+2)$.
\end{exercici}

\begin{solucio}
Si dividim $p(x)$ per $(x+1)(x+2)$, segons l'algorisme de la divisió, obtindrem un quocient $c(x)$ i un residu $r(x)$ de manera que\[
p(x)=(x+1)(x+2)\cdot c(x)+r(x)
\]i $\operatorname{grau}~r(x)<2$; d'aquesta darrera condició obtenim que $r(x)$ serà un polinomi de primer grau, o sigui de la forma $ax+b$. Utilitzant
ara les condicions de l'enunciat, pel teorema del residu, tenim\[
p(-1)=(-1+1)(-1+2)\cdot c(-1)+r(-1)=0
\]i\[
p(-2)=(-2+1)(-2+2)\cdot c(-2)+r(-2)=4
\]Aleshores, es compleix\begin{eqnarray*}
r(-1) &=&-a+b=0 \\
r(-2) &=&-2a+b=4
\end{eqnarray*}Resolent aquest sistema, obtenim que $a=b=-4$ i, per tant, el residu es\[
r(x)=-4x-4
\]
\end{solucio}

\begin{exercici}
Troba el valor de $a$ perquè el polinomi $p(x)=x^{4}-2x^{3}+ax^{2}-5x+10$
sigui divisible per $x+2$.
\end{exercici}

\begin{solucio}
Recorda que, segons el teorema del residu, el valor numèric del
polinomi $p(x)$ per a $x=a$ coincideix amb el residu de dividir $p(x)$ per $x-a$. Per tant, s'ha de verificar que\[
p(-2)=0
\]és a dir,\begin{eqnarray*}
\left( -2\right) ^{4}-2\left( -2\right) ^{3}+a\left( -2\right) ^{2}-5\left(
-2\right) +10 &=&0 \\
16+16+4a+10+10 &=&0 \\
4a &=&-52 \\
a &=&-13
\end{eqnarray*}
\end{solucio}

\begin{exercici}
Troba $a$ i $b$ perquè el polinomi $p(x)=x^{2}+ax-b$ doni de residu $1$
en dividir-lo per $x-1$, i residu $3$ en dividir-lo per $x-3$.
\end{exercici}

\begin{solucio}
Recorda que, segons el teorema del residu, el valor numèric del
polinomi $p(x)$ per a $x=a$ coincideix amb el residu de dividir $p(x)$ per $x-a$. Per tant, s'ha de verificar que\[
p(1)=1
\]i\[
p(3)=3
\]és a dir, s'han de satisfer les dues equacions següents\begin{eqnarray*}
1^{2}+a\cdot 1-b &=&1 \\
3^{2}+a\cdot 3-b &=&3
\end{eqnarray*}Resolent aquest sistema, obtenim $a=b=-3$.
\end{solucio}

\section{Arrels i descomposició factorial d'un polinomi}

\begin{exercici}
Comprova que els valors indicades són arrels dels polinomis
que s'acompanyen\[
\begin{array}{lll}
a) & 1 & p(x)=x^{3}-1 \\
b) & 1/2 & q(x)=16x^{3}-2x^{2}+x-2 \\
c) & \sqrt{3} & r(x)=x^{4}-5x^{2}+6\end{array}\]
\end{exercici}

\begin{solucio}
Recorda que diem que un nombre real $a$ és una arrel d'un
polinomi $p(x)$ si $p(a)=0$.

a) Tenim\[
p(1)=1^{3}-1=0
\]després, $1$ és arrel del polinomi.

b) Tenim\begin{eqnarray*}
q\left( \frac{1}{2}\right) &=&16\left( \frac{1}{2}\right) ^{3}-2\left( \frac{1}{2}\right) ^{2}+\frac{1}{2}-2 \\
&=&2-\frac{1}{2}+\frac{1}{2}-2 \\
&=&0
\end{eqnarray*}després, $1/2$ és arrel del polinomi.

c) Tenim\begin{eqnarray*}
r\left( \sqrt{3}\right) &=&\left( \sqrt{3}\right) ^{4}-5\left( \sqrt{3}\right) ^{2}+6 \\
&=&9-15+6 \\
&=&0
\end{eqnarray*}després, $\sqrt{3}$ és arrel del polinomi.
\end{solucio}

\begin{exercici}
Comprova si són certes o falses les afirmacions següents: (a) $1$ és
una arrel doble del polinomi $x^{3}+1$, (b) $\frac{1+\sqrt{2}}{2}$ és
una arrel simple del polinomi $4x^{2}-4x-1$, i (c) $2$ és una arrel
triple del polinomi $x^{4}-6x^{3}+12x^{2}-8x$.
\end{exercici}

\begin{solucio}
a) Observem que $1$ no és arrel del polinomi $p(x)=x^{3}+1$, ja que $p(1)\neq 0$, i, per tant, no és arrel doble.

b) Si fem $p(x)=4x^{2}-4x-1$, aleshores tenim\begin{eqnarray*}
p\left( \frac{1+\sqrt{2}}{2}\right) &=&4\left( \frac{1+\sqrt{2}}{2}\right)
^{2}-4\left( \frac{1+\sqrt{2}}{2}\right) -1 \\
&=&1+2\sqrt{2}+2-2-2\sqrt{2}-1 \\
&=&0
\end{eqnarray*}després, $\frac{1+\sqrt{2}}{2}$ és arrel de $p(x)$. Si ara dividim $p(x) $per $x-\frac{1+\sqrt{2}}{2}$, obtenim\[
4x^{2}-4x-1=\left( x-\frac{1+\sqrt{2}}{2}\right) \left( 4x-2+2\sqrt{2}\right)
\]i si fem $c(x)=4x-2+2\sqrt{2}$, aleshores\begin{eqnarray*}
c\left( \frac{1+\sqrt{2}}{2}\right) &=&4\left( \frac{1+\sqrt{2}}{2}\right)
-2+2\sqrt{2} \\
&=&2+2\sqrt{2}-2+2\sqrt{2} \\
&\neq &0
\end{eqnarray*}i, per tant, $\frac{1+\sqrt{2}}{2}$ és arrel simple de $p(x)$.

c) Si fem $p(x)=x^{4}-6x^{3}+12x^{2}-8x$, aleshores tenim\begin{eqnarray*}
p(2) &=&2^{4}-6\cdot 2^{3}+12\cdot 2^{2}-8\cdot 2 \\
&=&16-48+48-16 \\
&=&0
\end{eqnarray*}i, per tant, $2$ és arrel de $p(x)$. Observem que\[
p(x)=x\cdot \left( x^{3}-6x^{2}+12x-8\right)
\]i ara, mitjançant la regla de Ruffini, tenim\[
\begin{array}{r|rrrr}
& 1 & -6 & 12 & -8 \\
2 & & 2 & -8 & 8 \\\hline
& 1 & -4 & 4 & 0 \\
2 & & 2 & -4 & \\\hline
& 1 & -2 & 0 & \\
2 & & 2 & & \\\hline
 & 1 & 0 & &
\end{array}\]i, per tant, deducimos que $2$ és arrel triple del polinomi.
\end{solucio}

\begin{exercici}
Troba la multiplicitat de l'arrel $2$ pel polinomi $p(x)=x^{4}-8x^{3}+24x^{2}-32x+16$.
\end{exercici}

\begin{solucio}
Mitjançant la regla de Ruffini, obtenim\[
\begin{array}{r|rrrrr}
& 1 & -8 & 24 & -32 & 16 \\
2 & & 2 & -12 & 24 & -16 \\\hline
& 1 & -6 & 12 & -8 & 0 \\
2 & & 2 & -8 & 8 & \\\hline
& 1 & -4 & 4 & 0 & \\
2 & & 2 & -4 & & \\\hline
& 1 & -2 & 0 & & \\
2 & & 2 & & & \\\hline
 & 1 & 0 & & &
\end{array}\]i, per tant, $2$ és una arrel múltiple d'ordre $4$. Observa que es
compleix\[
x^{4}-8x^{3}+24x^{2}-32x+16=(x-2)^{4}
\]
\end{solucio}

\begin{exercici}
Calcula les arrels enteres del polinomi $p(x)=x^{5}-2x^{4}-6x^{3}-6x^{2}-7x-4$. Quines són les seves
multiplicitats? Té $p(x)$ alguna arrel real que no
sigui entera?
\end{exercici}

\begin{solucio}
Com que el polinomi $p(x)$ té tots els seus coeficients enters, les arrels enteres s'obtindran dels divisors del seu terme
independent. Aleshores, les possibles arrels enteres són $\pm 1,\pm 2$ o $\pm 4$. Mitjançant la regla de Ruffini podem ara esbrinar quines de
són efectivament arrels del polinomi.\[
\begin{array}{r|rrrrrr}
& 1 & -2 & -6 & -6 & -7 & -4 \\
-1 & & -1 & 3 & 3 & 3 & 4 \\\hline
& 1 & -3 & -3 & -3 & -4 & 0 \\
-1 & & -1 & 4 & -1 & 4 & \\\hline
& 1 & -4 & 1 & -4 & 0 & \\
4 & & 4 & 0 & 4 & & \\\hline
 & 1 & 0 & 1 & & &
\end{array}\]Per tant, les arrels enteres de $p(x)$ són $-1$ i $4$; a més, $-1$
és arrel doble de $p(x)$. D'aquesta manera, hem obtingut la següent
descomposició factorial del polinomi\[
p(x)=(x+1)^{2}\cdot (x-4)\cdot (x^{2}+1)
\]Observa que $x^{2}+1$ no té arrels, perquè no hi ha cap nombre real que el seu quadrat sumat amb $1$ doni com a resultat zero. Aleshores, $p(x)$ no té cap arrel real que no sigui entera.
\end{solucio}

\begin{exercici}
Escriu un polinomi de tercer grau que tingui $-2$ com a arrel múltiple d'ordre dos.
\end{exercici}

\begin{solucio}
Si el polinomi $p(x)$ ha de tenir a $-2$ com a arrel de multiplicitat $2$, el seu descomposició factorial contindrà el factor $(x+2)^{2}$. Si
ara fem que també tingui, per exemple, l'arrel $0$, aleshores\begin{eqnarray*}
p(x) &=&(x+2)^{2}\cdot x \\
&=&x^{3}+4x^{2}+4x
\end{eqnarray*}Per tant, $x^{3}+4x^{2}+4x$ és un polinomi de tercer grau que té $-2$
com a arrel doble.
\end{solucio}

\begin{exercici}
Calcula el valor que ha de tenir $a$ perquè el polinomi $p(x)=x^{3}-3x+a$
tingui l'arrel $1$.
\end{exercici}

\begin{solucio}
Recorda que diem que un nombre real $a$ és una arrel d'un
polinomi $p(x)$ si $p(a)=0$. Aleshores, s'ha de complir que\[
p(1)=0
\]és a dir,\begin{eqnarray*}
1^{3}-3\cdot 1+a &=&0 \\
a &=&2
\end{eqnarray*}
\end{solucio}

\begin{exercici}
Indica les arrels del polinomi $p(x)=(x-1)(x+\sqrt{2})(x-\sqrt{2})$.
\end{exercici}

\begin{solucio}
Com el polinomi $p(x)$ ve descompost en factors lineals, obtenim
de seguida els seus arrels. Les seves arrels són $1$, $-\sqrt{2}$ i $\sqrt{2}$.
\end{solucio}

\begin{exercici}
Troba un polinomi de segon grau que tingui per arrels $\sqrt{3}$ i $-\sqrt{2}$.
\end{exercici}

\begin{solucio}
Si el polinomi $p(x)$ que busquem és de segon grau i té arrels $\sqrt{3}$ i $-\sqrt{2}$, aleshores $p(x)$ admetirá la següent
descomposició en factors lineals\begin{eqnarray*}
p(x) &=&\left( x-\sqrt{3}\right) \cdot \left( x+\sqrt{2}\right) \\
&=&x^{2}+\left( \sqrt{2}-\sqrt{3}\right) x-\sqrt{6}
\end{eqnarray*}
\end{solucio}

\begin{exercici}
Calcula la descomposició factorial dels polinomis següents:\[
\begin{array}{l}
a)~(x-1)^{2}-(x-1)(x+1) \\
b)~(x-1)(x^{2}-4) \\
c)~x^{4}+2x^{3}-2x-1 \\
d)~x^{3}-1 \\
e)~3x^{4}-27x^{3}+39x^{2}+27x-42\end{array}\]
\end{exercici}

\begin{solucio}
a) Tenim,\begin{eqnarray*}
(x-1)^{2}-(x-1)\cdot (x+1) &=&(x-1)\cdot \left( x-1-x-1\right) \\
&=&-2\cdot (x-1)
\end{eqnarray*}en on hem sacado factor comú per a obtenir la descomposició.

b) Tenim,\[
(x-1)(x^{2}-4)=(x-1)\cdot (x-2)\cdot (x+2)
\]perquè, sabem que $x^{2}-4=(x-2)\cdot (x+2)$.

c) Mitjançant la regla de Ruffini, obtenim\[
\begin{array}{r|rrrrr}
& 1 & 2 & 0 & -2 & -1 \\
1 & & 1 & 3 & 3 & 1 \\\hline
& 1 & 3 & 3 & 1 & 0 \\
-1 & & -1 & -2 & -1 & \\\hline
& 1 & 2 & 1 & 0 & \\
-1 & & -1 & -1 & & \\\hline
& 1 & 1 & 0 & & \\
-1 & & -1 & & & \\\hline
 & 1 & 0 & & &
\end{array}\]Per tant, tenim la següent descomposició factorial\[
x^{4}+2x^{3}-2x-1=(x-1)\cdot (x+1)^{3}
\]

d) Mitjançant la regla de Ruffini, obtenim\[
\begin{array}{r|rrrr}
& 1 & 0 & 0 & -1 \\
1 & & 1 & 1 & 1 \\\hline
& 1 & 1 & 1 & 0\end{array}\]Per tant, tenim\[
x^{3}-1=(x-1)\cdot (x^{2}+x+1)
\]Ara bé $x^{2}+x+1$ no admet arrels reals, perquè, si no fos així,
l'equació de segon grau següent\[
x^{2}+x+1=0
\]tindria solucions i això no és possible perquè\[
b^{2}-4ac=1-4<0
\]En resum, la descomposició factorial del polinomi és\[
x^{3}-1=(x-1)\cdot (x^{2}+x+1)
\]

e) Observa primer que es compleix\[
3x^{4}-27x^{3}+39x^{2}+27x-42=3\cdot \left(
x^{4}-9x^{3}+13x^{2}+9x-14\right)
\]Ara, mitjançant la regla de Ruffini, obtenim\[
\begin{array}{r|rrrrr}
& 1 & -9 & 13 & 9 & -14 \\
1 & & 1 & -8 & 5 & 14 \\\hline
& 1 & -8 & 5 & 14 & 0 \\
-1 & & -1 & 9 & -14 & \\\hline
& 1 & -9 & 14 & 0 & \\
2 & & 2 & -14 & & \\\hline
& 1 & -7 & 0 & & \\
7 & & 7 & & & \\\hline
 & 1 & 0 & & &
\end{array}\]Per tant, tenim\[
3x^{4}-27x^{3}+39x^{2}+27x-42=3\cdot (x-1)\cdot (x+1)\cdot (x-2)\cdot (x-7)
\]
\end{solucio}

\begin{exercici}
Troba el m.c.d. i m.c.m. de $p(x)$ i $q(x)$ en els casos següents:\[
\begin{array}{l}
a)~p(x)=(x-1)^{2}(x-2)\text{ i }q(x)=(x+1)(x+2)^{2} \\
b)~p(x)=(x-\sqrt{2})^{2}(x+3)\text{ i }q(x)=(x-\sqrt{2})^{3}(x+\sqrt{2}) \\
c)~p(x)=x^{6}-1\text{ i }q(x)=x^{3}-x \\
d)~p(x)=3x^{4}-14x^{3}+20x^{2}-11x+2\text{ i }q(x)=3x^{4}-5x^{3}-4x^{2}+10x-4\end{array}\]
\end{exercici}

\begin{solucio}
Recorda que el màxim comú divisor s'obtindrà multiplicant
els factors irreduïbles comuns elevats a l'exponent més petit
que figuri en les dues descomposicions, i el mínim comú múltiple s'obtindrà multiplicant els factors irreduïbles comuns i no
comuns elevats a l'exponent més gran que figuri en totes dues
descomposicions.

a) Com que els dos polinomis estan donats pels seus descomposicions
factorials, tenim de seguida que\[
\text{m.c.d.}\left( p(x),q(x)\right) =1
\]i\[
\text{m.c.m.}\left( p(x),q(x)\right) =(x-1)^{2}\cdot (x-2)\cdot (x+1)\cdot
(x+2)^{2}
\]

b) Per la mateixa raó, tenim que\[
\text{m.c.d.}\left( p(x),q(x)\right) =(x-\sqrt{2})^{2}
\]i\[
\text{m.c.m.}\left( p(x),q(x)\right) =(x-\sqrt{2})^{3}\cdot (x+3)\cdot (x+\sqrt{2})
\]

c) En aquest cas, cal primer calcular les descomposicions dels dos
polinomis. Per la regla de Ruffini, obtenim\[
\begin{array}{r|rrrrrrr}
& 1 & 0 & 0 & 0 & 0 & 0 & -1 \\
1 & & 1 & 1 & 1 & 1 & 1 & 1 \\\hline
& 1 & 1 & 1 & 1 & 1 & 1 & 0 \\
-1 & & -1 & 0 & -1 & 0 & -1 & \\\hline
 & 1 & 0 & 1 & 0 & 1 & 0 &
\end{array}\]i, per tant,
\[
x^{6}-1=(x-1)\cdot (x+1)\cdot (x^{4}+x^{2}+1)
\]Observa que el factor $x^{4}+x^{2}+1$ és irreduïble, perquè, si no ho fos,
l'equació (bicuadrada)\[
x^{4}+x^{2}+1=0
\]admetiría solucióes, però, si fem el cambio $x^{2}=t$, aleshores
obtenim l'equació de segon grau\[
t^{2}+t+1=0
\]que no té solucions reals, perquè,\[
b^{2}-4ac=1-4<0
\]D'altra banda, tenim\[
x^{3}-x=x\cdot (x^{2}-1)=x\cdot (x-1)\cdot (x+1)
\]Per consegüent, tenim\[
\begin{array}{l}
p(x)=(x-1)\cdot (x+1)\cdot (x^{4}+x^{2}+1) \\
q(x)=x\cdot (x-1)\cdot (x+1)\end{array}\]i, per tant,\[
\text{m.c.d.}\left( p(x),q(x)\right) =(x-1)\cdot (x+1)
\]i\[
\text{m.c.m.}\left( p(x),q(x)\right) =x\cdot (x-1)\cdot (x+1)\cdot
(x^{4}+x^{2}+1)
\]

d) Calcularem primer la descomposició factorial de $p(x)=3x^{4}-14x^{3}+20x^{2}-11x+2$. Com tots els seus coeficients són nombres enters, les arrels enteres d'aquest polinomi s'obtindran a
partir dels divisors del terme independent. En aquest cas, els
possibles candidats són $\pm 1$ i $\pm 2$. Mitjançant la regla de Ruffini,
observem que només $1$ és arrel entera de $p(x)$, perquè,\[
\begin{array}{r|rrrrr}
& 3 & -14 & 20 & -11 & 2 \\
1 & & 3 & -11 & 9 & -2 \\\hline
& 3 & -11 & 9 & -2 & 0\end{array}\]D'aquí, obtenim\[
p(x)=(x-1)\cdot (3x^{3}-11x^{2}+9x-2)
\]Les altres possibles arrels, pertanyents el polinomi $3x^{3}-11x^{2}+9x-2$, seran nombres reals que no són nombres
enters. Com que calcular-les directament és feixuc i encara tenim un polinomi de tercer grau, veiem que per a calcular el m.c.d. de aquests dos polinomis
serà més convenient usar l'algorisme d'Euclides. Així, tenim\[
\begin{array}{ccccc|ccccc|cccc|ccc}
& & & & & 1 & & & & & x & +1 & & & 3x & -3 &
\\\cline{6-17} 3x^{4} & -14x^{3} & +20x^{2} & -11x & +2 & 3x^{4} &
-5x^{3} & -4x^{2} & +10x
& -4 & 3x^{3} & -8x^{2} & +7x & -2 & x^{2} & -\frac{5}{3}x & +\frac{2}{3} \\
-3x^{4} & +5x^{3} & +4x^{2} & -10x & +4 & -3x^{4} & +8x^{3} &
-7x^{2} & +2x & & -3x^{3} & +5x^{2} & -2x & & & &
\\\cline{1-5} \cline{6-9} \cline{11-13} \diagup & -9x^{3} & +24x^{2} & -21x & +6 & \diagup &
3x^{3} & -11x^{2} & 12x
& -4 & \diagup & -3x^{2} & +5x & -2 & & & \\
& 3x^{3} & -8x^{2} & +7x & -2 & & -3x^{3} & +8x^{2} & -7x & +2 &
& +3x^{2} & -5x & +2 & & & \\\cline{7-10} \cline{12-14} & & &
& & & \diagup & -3x^{2} & +5x & -2 & & \diagup & 0 & 0 & & &
\\
& & & & & & & x^{2} & -\frac{5}{3}x & +\frac{2}{3} & & & &
& & &
\end{array}\]Per tant, tenim\begin{eqnarray*}
\text{m.c.d.}\left( p(x),q(x)\right) &=&\text{m.c.d.}\left(
3x^{4}-5x^{3}-4x^{2}+10x-4,3x^{3}-8x^{2}+7x-2\right) \\
&=&\text{m.c.d.}\left( 3x^{3}-8x^{2}+7x-2,x^{2}-\frac{5}{3}x+\frac{2}{3}\right) \\
&=&\text{m.c.d.}\left( x^{2}-\frac{5}{3}x+\frac{2}{3},0\right) \\
&=&x^{2}-\frac{5}{3}x+\frac{2}{3}
\end{eqnarray*}i, aleshores, obtenim\begin{eqnarray*}
\text{m.c.m.}\left( p(x),q(x)\right) &=&\frac{1}{3\cdot 3}\cdot \frac{\left(
3x^{4}-14x^{3}+20x^{2}-11x+2\right) \cdot \left(
3x^{4}-5x^{3}-4x^{2}+10x-4\right) }{x^{2}-\frac{5}{3}x+\frac{2}{3}} \\
&=&x^{6}-\frac{14}{3}x^{5}+\frac{14}{3}x^{4}+\frac{17}{3}x^{3}-\frac{38}{3}x^{2}+\frac{22}{3}x-\frac{4}{3}
\end{eqnarray*}
\end{solucio}

\begin{exercici}
Dos polinomis s'anomenen primers entre si quan el seu m.c.d. és $1$, és a
dir, quan el únic divisor comú és el polinomi unitat. Indica
quines de les parelles següents de polinomis són primers entre si\[
\begin{array}{l}
a)~p(x)=x^{2}-2x+1\text{ i }q(x)=x^{2}+2x+1 \\
b)~p(x)=x^{3}-3x^{2}\text{ i }q(x)=x^{3}+3x^{2} \\
c)~p(x)=x^{2}+x+1\text{ i }q(x)=x^{2}-x-1\end{array}\]
\end{exercici}

\begin{solucio}
a) Com que\[
x^{2}-2x+1=(x-1)^{2}
\]i\[
x^{2}+2x+1=(x+1)^{2}
\]aleshores,\[
\text{m.c.d.}\left( p(x),q(x)\right) =1
\]i els dos polinomis són primers entre si.

b) Com que\[
x^{3}-3x^{2}=x^{2}\cdot (x-3)
\]i\[
x^{3}+3x^{2}=x^{2}\cdot (x+3)
\]aleshores,\[
\text{m.c.d.}\left( p(x),q(x)\right) =x^{2}
\]i els dos polinomis no són primers entre si.

c) Com que $p(x)$ és irreduïble i $\operatorname{grau}~q(x)=2$, és clar que els dos
polinomis són primers entre si.
\end{solucio}

\begin{exercici}
Resol les següents equacions polinòmiques\[
\begin{array}{l}
a)~x^{3}+21x^{2}+98x-120=0 \\
b)~x^{3}-14x^{2}+63x-90=0\end{array}\]
\end{exercici}

\begin{solucio}
Observa que buscar les solucions reals de l'equació
\[
p(x)=0
\]és equivalent, per definició, a trobar les arrels del
polinomi $p(x)$.

a) Com tots els coeficients del polinomi són nombres enters, les arrels enteres es trobaran entre els divisors del terme
independent. Provarem aquests divisors mitjançant la regla de Ruffini,\[
\begin{array}{r|rrrr}
& 1 & 21 & 98 & -120 \\
1 & & 1 & 22 & 120 \\\hline
& 1 & 22 & 120 & 0\end{array}\]D'aquí, obtenim\[
x^{3}+21x^{2}+98x-120=(x-1)\cdot (x^{2}+22x+120)
\]Podríem continuar buscant més arrels per Ruffini, però serà més senzill procedir d'una altra manera. Observa que l'equació
inicial pot ara escriure's de la manera següent\[
(x-1)\cdot (x^{2}+22x+120)=0
\]i, per tant, les solucions de l'equació són\[
\begin{array}{ccc}
x-1=0 & \Longrightarrow & x=1\end{array}\]i les altres solucions s'obtindran de l'equació següent de
segon grau
\[
x^{2}+22x+120=0
\]que, resolviéndola, obtenim $x=-10$ o $x=-12$. Per consegüent, les
solucions de l'equació
\[
x^{3}+21x^{2}+98x-120=0
\]són $1$, $-10$ i $-12$.

b) Procederemos de la mateixa manera, primer, mitjançant la regla de Ruffini,
obtenim\[
\begin{array}{r|rrrr}
& 1 & -14 & 63 & -90 \\
3 & & 3 & -33 & 90 \\\hline
& 1 & -11 & 30 & 0\end{array}\]Per tant, l'equació original pot escriure's de la manera següent\[
(x-3)\cdot (x^{2}-11x+30)=0
\]Les solucions d'aquesta equació són\[
\begin{array}{ccc}
x-3=0 & \Longrightarrow & x=3\end{array}\]i les que s'obtenen de l'equació següent de segon grau\[
x^{2}-11x+30=0
\]les quals solucions són $5$ i $6$. Per consegüent, les solucions de l'equació
\[
x^{3}-14x^{2}+63x-90=0
\]són $3$, $5$ i $6$.
\end{solucio}

\section{El cos de les fraccions racionals}

\begin{exercici}
Indica si són o no equivalents les parelles següents de fraccions
algebraiques:\[
\begin{array}{cccc}
a) & \frac{x^{2}-1}{x^{3}-1} & & \frac{x+1}{x^{2}+x+1} \\
b) & \frac{x^{3}+x^{2}+x+1}{x^{2}+1} & & \frac{x^{2}+x}{x} \\
c) & \frac{1}{x} & & \frac{x+1}{x^{2}+x}\end{array}\]
\end{exercici}

\begin{solucio}
Recorda que diem que dues fraccions algebraiques\[
\frac{p(x)}{q(x)}\text{ \ \ \ i \ \ \ }\frac{r(x)}{s(x)}
\]són equivalents si i només si\[
p(x)\cdot s(x)=q(x)\cdot r(x)
\]

a) Com que\[
\left( x^{2}-1\right) \cdot \left( x^{2}+x+1\right) =x^{4}+x^{3}-x-1
\]i\[
\left( x^{3}-1\right) \cdot \left( x+1\right) =x^{4}+x^{3}-x-1
\]es té
\[
\frac{x^{2}-1}{x^{3}-1}=\frac{x+1}{x^{2}+x+1}
\]o sigui, són equivalents.

b) Com que\[
\left( x^{3}+x^{2}+x+1\right) \cdot x=x^{4}+x^{4}+x^{2}+x
\]i\[
\left( x^{2}+1\right) \cdot \left( x^{2}+x\right) =x^{4}+x^{3}+x^{2}+x
\]es té\[
\frac{x^{3}+x^{2}+x+1}{x^{2}+1}=\frac{x^{2}+x}{x}
\]

c) Com que\[
1\cdot \left( x^{2}+x\right) =x\cdot (x+1)
\]també es té\[
\frac{1}{x}=\frac{x+1}{x^{2}+x}
\]
\end{solucio}

\begin{exercici}
Escriu tres fraccions algebraiques pertanyents a la fracció
racional\[
\left\{ \frac{x}{x+3}\right\}
\]
\end{exercici}

\begin{solucio}
De la definició d'equivalència de fraccions algebraiques, és clar que
es compleix\begin{eqnarray*}
\frac{x}{x+3} &=&\frac{x^{2}}{x^{2}+3x} \\
&=&\frac{x^{2}-x}{x^{2}+2x-3} \\
&=&\frac{x^{3}-x^{2}}{x^{3}+2x^{2}-3x}
\end{eqnarray*}
\end{solucio}

\begin{exercici}
Simplificar les fraccions algebraiques següents:\[
\begin{array}{lllll}
a)~\frac{x+2}{x^{2}-4} & & b)~\frac{x-\sqrt{2}}{x^{2}-2\sqrt{2}x+2} & & c)~\frac{x}{x^{4}-x} \\
d)~\frac{2x^{2}-3x+1}{3x^{3}-3x^{2}+x-1} & & e)~\frac{(x+3)^{2}-(5x-4)^{2}}{36x^{2}-1} & & f)~\frac{x^{2}-1}{x^{3}-2x^{2}+x}\end{array}\]
\end{exercici}

\begin{solucio}
Recorda que simplificar una fracció algebraica és buscar la fracció irreduïble equivalent.

a)\begin{eqnarray*}
\frac{x+2}{x^{2}-4} &=&\frac{x+2}{(x+2)(x-2)} \\
&=&\frac{1}{x-2}
\end{eqnarray*}

b)\begin{eqnarray*}
\frac{x-\sqrt{2}}{x^{2}-2\sqrt{2}x+2} &=&\frac{x-\sqrt{2}}{\left( x-\sqrt{2}\right) ^{2}} \\
&=&\frac{1}{x-\sqrt{2}}
\end{eqnarray*}

c)\begin{eqnarray*}
\frac{x}{x^{4}-x} &=&\frac{x}{x(x^{3}-1)} \\
&=&\frac{1}{x^{3}-1}
\end{eqnarray*}

d)\begin{eqnarray*}
\frac{2x^{2}-3x+1}{3x^{3}-3x^{2}+x-1} &=&\frac{(x-1)(2x-1)}{(x-1)(3x^{2}+1)}
\\
&=&\frac{2x-1}{3x^{2}+1}
\end{eqnarray*}perquè, per la regla de Ruffini, tenim\[
\begin{array}{r|rrr}
& 2 & -3 & 1 \\
1 & & 2 & -1 \\\hline
& 2 & -1 & 0\end{array}\]i\[
\begin{array}{r|rrrr}
& 3 & -3 & 1 & -1 \\
1 & & 3 & 0 & 1 \\\hline
& 3 & 0 & 1 & 0\end{array}\]

e)\begin{eqnarray*}
\frac{(x+3)^{2}-(5x-4)^{2}}{36x^{2}-1} &=&\frac{\left( x+3+5x-4\right) \cdot
\left( x+3-5x+4\right) }{(6x+1)(6x-1)} \\
&=&\frac{(6x-1)(-4x+7)}{(6x+1)(6x-1)} \\
&=&\frac{-4x+7}{6x+1}
\end{eqnarray*}

f)\begin{eqnarray*}
\frac{x^{2}-1}{x^{3}-2x^{2}+x} &=&\frac{(x+1)(x-1)}{x(x^{2}-2x+1)} \\
&=&\frac{(x+1)(x-1)}{x(x-1)^{2}} \\
&=&\frac{x+1}{x(x-1)}
\end{eqnarray*}
\end{solucio}

\begin{exercici}
Efectua les operacions indicades\[
\begin{array}{ll}
a) & \frac{1+x}{1-x}+x-\frac{x}{1-x^{2}} \\
b) & 3-\frac{1}{1+x}+\frac{x^{2}-1}{1-x} \\
c) & \frac{x+1}{3}+\frac{x^{2}-1}{9(x+1)}+\frac{x+1}{6(x^{2}-1)} \\
d) & \frac{1}{x-1}+\frac{1}{x^{2}-1}-\frac{1}{x+1}\end{array}\]
\end{exercici}

\begin{solucio}
a)\begin{eqnarray*}
\frac{1+x}{1-x}+x-\frac{x}{1-x^{2}} &=&\frac{(1+x)^{2}}{1-x^{2}}+\frac{x(1-x^{2})}{1-x^{2}}-\frac{x}{1-x^{2}} \\
&=&\frac{1+2x+x^{2}+x-x^{3}-x}{1-x^{2}} \\
&=&\frac{1+2x+x^{2}-x^{3}}{1-x^{2}}
\end{eqnarray*}

b)\begin{eqnarray*}
3-\frac{1}{1+x}+\frac{x^{2}-1}{1-x} &=&\frac{3(1+x)(1-x)}{(1+x)(1-x)}-\frac{1-x}{(1+x)(1-x)}+\frac{(1+x)(x^{2}-1)}{(1+x)(1-x)} \\
&=&\frac{3-3x^{2}-1+x+x^{2}-1+x^{3}-x}{(1+x)(1-x)} \\
&=&\frac{x^{3}-2x^{2}+1}{(1+x)(1-x)} \\
&=&-\frac{(1-x)(x^{2}-x-1)}{(1+x)(1-x)} \\
&=&-\frac{x^{2}-x-1}{1+x}
\end{eqnarray*}

c)\begin{eqnarray*}
\frac{x+1}{3}+\frac{x^{2}-1}{9(x+1)}+\frac{x+1}{6(x^{2}-1)} &=&\frac{6(x+1)^{2}(x-1)}{18(x^{2}-1)}+\frac{2(x-1)(x^{2}-1)}{18(x^{2}-1)}+\frac{3(x+1)}{18(x^{2}-1)} \\
&=&\frac{-6+6x^{2}-6x+6x^{3}+2-2x^{2}-2x+2x^{3}+3x+3}{18(x^{2}-1)} \\
&=&\frac{8x^{3}+4x^{2}-5x-1}{18(x^{2}-1)} \\
&=&\frac{(x+1)(8x^{2}-4x-1)}{18(x+1)(x-1)} \\
&=&\frac{8x^{2}-4x-1}{18(x-1)}
\end{eqnarray*}

d)\begin{eqnarray*}
\frac{1}{x-1}+\frac{1}{x^{2}-1}-\frac{1}{x+1} &=&\frac{x+1}{x^{2}-1}+\frac{1}{x^{2}-1}-\frac{x-1}{x^{2}-1} \\
&=&\frac{x+1+1-x+1}{x^{2}-1} \\
&=&\frac{3}{x^{2}-1}
\end{eqnarray*}
\end{solucio}

\begin{exercici}
Efectua les operacions, simplificant al màxim el
resultat:\[
\begin{array}{ll}
a) & \frac{1+x}{1-x}\cdot \frac{x}{1-x^{2}} \\
b) & \frac{3x^{2}+1}{x^{2}+2x+1}\cdot \frac{x^{2}-1}{x+2} \\
c) & \frac{x^{2}-4x+4}{x+1}\cdot \frac{x^{2}-1}{x-2} \\
d) & \frac{8x}{x+1}\cdot \frac{3x}{x-1}\cdot \frac{x^{2}-1}{x^{2}+1}\end{array}\]
\end{exercici}

\begin{solucio}
a)\begin{eqnarray*}
\frac{1+x}{1-x}\cdot \frac{x}{1-x^{2}} &=&\frac{x(1+x)}{(1-x)^{2}(1+x)} \\
&=&\frac{x}{(1-x)^{2}} \\
&=&\frac{x}{1-2x+x^{2}}
\end{eqnarray*}

b)\begin{eqnarray*}
\frac{3x^{2}+1}{x^{2}+2x+1}\cdot \frac{x^{2}-1}{x+2} &=&\frac{(3x^{2}+1)(x+1)(x-1)}{(x+1)^{2}(x+2)} \\
&=&\frac{(3x^{2}+1)(x-1)}{(x+1)(x+2)} \\
&=&\frac{3x^{3}-3x^{2}+x-1}{x^{2}+3x+2}
\end{eqnarray*}

c)\begin{eqnarray*}
\frac{x^{2}-4x+4}{x+1}\cdot \frac{x^{2}-1}{x-2} &=&\frac{(x-2)^{2}(x+1)(x-1)}{(x+1)(x-2)} \\
&=&\frac{(x-2)(x-1)}{1} \\
&=&x^{2}-3x+2
\end{eqnarray*}

d)\begin{eqnarray*}
\frac{8x}{x+1}\cdot \frac{3x}{x-1}\cdot \frac{x^{2}-1}{x^{2}+1} &=&\frac{8x}{x+1}\cdot \left( \frac{3x}{x-1}\cdot \frac{x^{2}-1}{x^{2}+1}\right) \\
&=&\frac{8x}{x+1}\cdot \frac{3x(x+1)(x-1)}{(x-1)(x^{2}+1)} \\
&=&\frac{8x}{x+1}\cdot \frac{3x(x+1)}{x^{2}+1} \\
&=&\frac{24x^{2}(x+1)}{(x+1)(x^{2}+1)} \\
&=&\frac{24x^{2}}{x^{2}+1}
\end{eqnarray*}
\end{solucio}

\begin{exercici}
Efectua les operacions, simplificant al màxim el
resultat:\[
\begin{array}{ll}
a) & \frac{x^{2}}{x+1}:\frac{3x^{2}+1}{x^{2}+2x+1} \\
b) & \left( \frac{1}{x}+\frac{1}{x+1}\right):\frac{x^{2}+x}{x^{3}-1} \\
c) & \frac{1}{x}\cdot \left( \frac{x^{2}}{x+1}-\frac{x^{3}}{x-1}\right) \\
d) & \frac{a(x+1)}{b}:\left( \frac{a+b}{x+1}-\frac{a-b}{bx+b}\right)\end{array}\]
\end{exercici}

\begin{solucio}
a)\begin{eqnarray*}
\frac{x^{2}}{x+1} &:&\frac{3x^{2}+1}{x^{2}+2x+1}=\frac{x^{2}}{x+1}\cdot
\frac{x^{2}+2x+1}{3x^{2}+1} \\
&=&\frac{x^{2}(x+1)^{2}}{(x+1)(3x^{2}+1)} \\
&=&\frac{x^{2}(x+1)}{3x^{2}+1} \\
&=&\frac{x^{3}+x^{2}}{3x^{2}+1}
\end{eqnarray*}

b)\begin{eqnarray*}
\left( \frac{1}{x}+\frac{1}{x+1}\right) &:&\frac{x^{2}+x}{x^{3}-1}=\left(
\frac{x+1}{x(x+1)}+\frac{x}{x(x+1)}\right):\frac{x^{2}+x}{x^{3}-1} \\
&=&\frac{x+1+x}{x(x+1)}:\frac{x^{2}+x}{x^{3}-1} \\
&=&\frac{2x+1}{x(x+1)}\cdot \frac{x^{3}-1}{x^{2}+x} \\
&=&\frac{(2x+1)(x^{3}-1)}{x^{2}(x+1)^{2}} \\
&=&\frac{2x^{4}+x^{3}-2x-1}{x^{4}+2x^{3}+x^{2}}
\end{eqnarray*}

c)\begin{eqnarray*}
\frac{1}{x}\cdot \left( \frac{x^{2}}{x+1}-\frac{x^{3}}{x-1}\right) &=&\frac{1}{x}\cdot \left[ \frac{x^{2}(x-1)}{(x+1)(x-1)}-\frac{x^{3}(x+1)}{(x+1)(x-1)}\right] \\
&=&\frac{1}{x}\cdot \frac{x^{2}(x-1-x^{2}-x)}{x^{2}-1} \\
&=&\frac{x^{2}(-1-x^{2})}{x(x^{2}-1)} \\
&=&-\frac{x+x^{3}}{x^{2}-1}
\end{eqnarray*}

d)\begin{eqnarray*}
\frac{a(x+1)}{b} &:&\left( \frac{a+b}{x+1}-\frac{a-b}{bx+b}\right) =\frac{a(x+1)}{b}:\left[ \frac{b(a+b)}{b(x+1)}-\frac{a-b}{b(x+1)}\right] \\
&=&\frac{a(x+1)}{b}:\frac{b(a+b)-a+b}{b(x+1)} \\
&=&\frac{a(x+1)}{b}\cdot \frac{b(x+1)}{b(a+b)-a+b} \\
&=&\frac{ab(x+1)^{2}}{b(ab+b^{2}-a+b)} \\
&=&\frac{a(x+1)^{2}}{ab+b^{2}-a+b}
\end{eqnarray*}
\end{solucio}

\begin{exercici}
Determinar la solució de l'equació\[
\frac{x^{2}}{x^{2}-1}\cdot A(x)=\frac{x^{2}-1}{x^{2}+x}+\frac{x+1}{x}
\]
\end{exercici}

\begin{solucio}
Resolent aquesta equació en el cos de les fraccions racionals,
tenim\begin{eqnarray*}
\frac{x^{2}-1}{x^{2}}\cdot \left( \frac{x^{2}}{x^{2}-1}\cdot A(x)\right) &=&\frac{x^{2}-1}{x^{2}}\cdot \left( \frac{x^{2}-1}{x^{2}+x}+\frac{x+1}{x}\right) \\
\left( \frac{x^{2}-1}{x^{2}}\cdot \frac{x^{2}}{x^{2}-1}\right) \cdot A(x) &=&\frac{x^{2}-1}{x^{2}}\cdot \left( \frac{x^{2}-1}{x^{2}+x}+\frac{x+1}{x}\right) \\
\frac{1}{1}\cdot A(x) &=&\frac{x^{2}-1}{x^{2}}\cdot \left[ \frac{x^{2}-1}{x(x+1)}+\frac{(x+1)^{2}}{x(x+1)}\right] \\
A(x) &=&\frac{x^{2}-1}{x^{2}}\cdot \frac{x^{2}-1+x^{2}+2x+1}{x(x+1)} \\
A(x) &=&\frac{x^{2}-1}{x^{2}}\cdot \frac{2x(x+1)}{x(x+1)} \\
A(x) &=&\frac{x^{2}-1}{x^{2}}\cdot \frac{2}{1} \\
A(x) &=&\frac{2x^{2}-2}{x^{2}}
\end{eqnarray*}
\end{solucio}

\begin{exercici}
Comprova que es compleix\[
\left( x+1+\frac{x^{2}}{1-x}\right):\left( 1-\frac{x}{1+x}\cdot \frac{x+1}{x^{2}-x}\right) =-\frac{1}{x-2}
\]
\end{exercici}

\begin{solucio}
\begin{eqnarray*}
\left( x+1+\frac{x^{2}}{1-x}\right) &:&\left( 1-\frac{x}{1+x}\cdot \frac{x+1}{x^{2}-x}\right) \\
&=&\left[ \frac{(x+1)(1-x)}{1-x}+\frac{x^{2}}{1-x}\right]:\left[ 1-\frac{x(x+1)}{x(x-1)(x+1)}\right] \\
&=&\frac{x-x^{2}+1-x+x^{2}}{1-x}:\left( 1-\frac{1}{x-1}\right) \\
&=&\frac{1}{1-x}:\frac{x-1-1}{x-1} \\
&=&\frac{1}{1-x}\cdot \frac{x-1}{x-2} \\
&=&-\frac{1-x}{(1-x)(x-2)} \\
&=&-\frac{1}{x-2}
\end{eqnarray*}
\end{solucio}

\begin{exercici}
Simplifica:

\begin{enumerate}
\item[a)]
\[
\frac{1+\left( \frac{1-x}{1+x}\right) ^{2}}{1-\left( \frac{1-x}{1+x}\right)
^{2}}
\]

\item[b)]
\[
\frac{1-\frac{x+a}{x-a}}{1+\frac{x+a}{x-a}}
\]
\end{enumerate}
\end{exercici}

\begin{solucio}
a)\begin{eqnarray*}
\frac{1+\left( \frac{1-x}{1+x}\right) ^{2}}{1-\left( \frac{1-x}{1+x}\right)
^{2}} &=&\left[ 1+\frac{(1-x)^{2}}{(1+x)^{2}}\right]:\left[ 1-\frac{(1-x)^{2}}{(1+x)^{2}}\right] \\
&=&\frac{1+2x+x^{2}+1-2x+x^{2}}{(1+x)^{2}}:\frac{1+2x+x^{2}-1+2x-x^{2}}{(1+x)^{2}} \\
&=&\frac{2(1+x^{2})}{(1+x)^{2}}\cdot \frac{(1+x)^{2}}{4x} \\
&=&\frac{2(1+x^{2})(1+x)^{2}}{4x(1+x)^{2}} \\
&=&\frac{1+x^{2}}{2x}
\end{eqnarray*}

b)\begin{eqnarray*}
\frac{1-\frac{x+a}{x-a}}{1+\frac{x+a}{x-a}} &=&\left( 1-\frac{x+a}{x-a}\right):\left( 1+\frac{x+a}{x-a}\right) \\
&=&\frac{x-a-x-a}{x-a}:\frac{x-a+x+a}{x-a} \\
&=&\frac{-2a}{x-a}\cdot \frac{x-a}{2x} \\
&=&\frac{-2a(x-a)}{2x(x-a)} \\
&=&-\frac{a}{x}
\end{eqnarray*}
\end{solucio}

\section{Funcions polinòmiques}

\begin{exercici}
Calcula $f(-1)$, $f(1/3)$ i $f(\sqrt{2})$ per a les funcions següents: (a)
$f(x)=2x+5$; (b) $f(x)=-x^{2}+2x-3$; i (c) $f(x)=1$.
\end{exercici}

\begin{solucio}
Es tracta de trobar les imatges dels punts $-1$, $1/3$ i $\sqrt{2}$
per cada una de les funcions donades.

a)\[
\begin{array}{l}
f(-1)=2\cdot (-1)+5=3 \\
f(1/3)=2\cdot \frac{1}{3}+5=\frac{17}{3} \\
f(\sqrt{2})=2\cdot \sqrt{2}+5\end{array}\]

b)\[
\begin{array}{l}
f(-1)=-(-1)^{2}+2\cdot (-1)-3=-1-2-3=-6 \\
f(1/3)=-\left( \frac{1}{3}\right) ^{2}+2\cdot \frac{1}{3}-3=-\frac{1}{9}+\frac{2}{3}-3=-\frac{22}{9} \\
f(\sqrt{2})=-\left( \sqrt{2}\right) ^{2}+2\cdot \sqrt{2}-3=-2+2\sqrt{2}-3=-5+2\sqrt{2}\end{array}\]

c)\[
\begin{array}{l}
f(-1)=1 \\
f(1/3)=1 \\
f(\sqrt{2})=1\end{array}\]perquè la funció és, en aquest cas, la funció constant igual a $1$,
és a dir, la imatge de qualsevol nombre real és sempre igual a $1$.
\end{solucio}

\begin{exercici}
Calcula, si existeixen, les antiimatges de $0$, $1$ i $-1$ per a les
funcions següents: (a) $f(x)=2x-1$; (b) $f(x)=2x^{2}$; (c) $f(x)=x^{3}$;
i (d) $f(x)=-1$.
\end{exercici}

\begin{solucio}
Recorda que donada una funció $f$, si $y$ és la imatge de $x$ per $f$,
escrivim\[
f(x)=y
\]i aleshores també es diu que $y$ és una antiimatge de $x$ per $f$. Es
tracta de trobar els valors de $x$ per als valors de $y$ donats en el
enunciat.

a)
\[
\begin{array}{lll}
f(x)=y & \Longrightarrow & x=~??? \\
2x-1=0 & \Longrightarrow & x=1/2 \\
2x-1=1 & \Longrightarrow & x=1 \\
2x-1=-1 & \Longrightarrow & x=0\end{array}\]

b)\[
\begin{array}{lll}
f(x)=y & \Longrightarrow & x=~??? \\
2x^{2}=0 & \Longrightarrow & x=0 \\
2x^{2}=1 & \Longrightarrow & x=\pm \sqrt{1/2} \\
2x^{2}=-1 & \Longrightarrow & \text{No hi ha antiimatges de }-1\end{array}\]

c)\[
\begin{array}{lll}
f(x)=y & \Longrightarrow & x=~??? \\
x^{3}=0 & \Longrightarrow & x=0 \\
x^{3}=1 & \Longrightarrow & x=1 \\
x^{3}=-1 & \Longrightarrow & x=\sqrt[3]{-1}=-1\end{array}\]

d) En aquest cas, la funció és constant igual a $-1$, és a dir, la
imatge de qualsevol nombre real per aquesta funció és sempre igual a $-1$. Per consegüent, només $-1$ té antiimatge per la funció i
aquesta és $-1$ i, qualsevol altre nombre real no té antiimatge.
\end{solucio}

\begin{exercici}
Construeix la gràfica de les funcions les imatges de les quals estan
definides per les relacions següents: (a) $y=-5$; (b) $y=4x$; (c) $y=-\frac{1}{2}x$; (d) $y=-3x+4$; (e) $y=3x+9$.
\end{exercici}

\begin{solucio}
Les gràfiques de totes les funcions indicades es mostran en la
següent figura. \begin{center}
\includegraphics[width=2.5901in,height=2.6524in,keepaspectratio]{poli15.png}
\end{center}

a) La relació $y=-5$ estableix una funció $f$ definida per $f(x)=-5 $, per a tot $x\in \mathbb{R}$. Es tracta d’una funció
constant i, per tant, la seva representació gràfica serà una recta
horizontal.

b) La relació $y=4x$ estableix una funció $f$ definida per $f(x)=4x $, per a tot $x\in \mathbb{R}$. Es tracta d’una funció lineal
i, per tant, la seva representació gràfica serà una recta de
pendent $4$ que passa per l’origen de coordenades. Com que el pendent
és positiva, la funció és creixent i pren valors positius per a $x>0$
i negatius, per a $x<0$.

c) La relació $y=-\frac{1}{2}x$ estableix una funció $f$ definida
per $f(x)=-\frac{1}{2}x$, per a tot $x\in \mathbb{R}$. Es tracta d’una funció lineal i, per tant, la seva representació gràfica serà una
recta de pendent $-\frac{1}{2}$ que passa per l’origen de coordenades.
Com que el pendent és negatiu, la funció és decreixent i pren
valors positius per a $x<0$ i negatius, per a $x>0$.

d) La relació $y=-3x+4$ estableix una funció $f$ definida per $f(x)=-3x+4$, per a tot $x\in \mathbb{R}$. Es tracta d’una funció afí i, per tant, la seva representació gràfica serà una recta
de pendent $-3$ que passa pel punt de coordenades $(0,4)$. Com que
el pendent és negatiu, la funció és decreixent i pren valors
positius per a $x<\frac{4}{3}$ i negatius, per a $x>\frac{4}{3}$, perquè la
funció s'anul·la per a $x=\frac{4}{3}$.

e) La relació $y=3x+9$ estableix una funció $f$ definida per $f(x)=3x+9$, per a tot $x\in \mathbb{R}$. Es tracta d’una funció afí i, per tant, la seva representació gràfica serà una recta de
pendent $3$ que passa pel punt de coordenades $(0,9)$. Com que la
pendent és positiu, la funció és creixent i pren valors positius
per a $x>-3$ i negatius, per a $x<-3$, perquè la funció s'anul·la per a $x=-3$.
\end{solucio}

\begin{exercici}
 En quins punts talla la gràfica de les funcions
indicades als eixos de coordenades? En quins punts
es tallen les dues gràfiques? Per a quins valors de $x$
les funcions prenen valors positius i valors negatius?\[
\begin{array}{l}
f(x)=-x+1 \\
g(x)=2x-4\end{array}\]
\end{exercici}

\begin{solucio}
Les gràfiques de les dues funcions indicades es mostran en la següent
figura. Evidentment, són funcions afins i, per tant, les gràfiques
són rectes. \begin{center}
\includegraphics[width=2.6109in,height=2.5901in,keepaspectratio]{poli16.png}
\end{center}

Els punts de tall amb l'eix d'abscisses s'obtenen com les antiimatges de $0$, és a dir, com els valors de $x$ que verifiquen\[
f(x)=0
\]i els punts de tall amb l'eix d’ordenades s'obtenen com les imatges de $0$. Així, per a la funció $f(x)=-x+1$, tenim\[
\begin{array}{ccc}
-x+1=0 & \Longrightarrow & x=1\end{array}\]i\[
f(0)=1
\]Per tant, la gràfica talla l'eix $X$ en el punt $(1,0)$ i l'eix $Y$
en $(0,1)$. Per a la funció $g(x)=2x-4$, tenim\[
\begin{array}{ccc}
2x-4=0 & \Longrightarrow & x=2\end{array}\]i\[
g(0)=-4
\]i, per tant, la gràfica talla l'eix $X$ en el punt $(2,0)$ i l'eix $Y$ en $(0,-4)$.

Els punts de tall de les dues gràfiques es corresponen amb els valors
de $x$ per als quals les imatges per totes dues funcions són iguals, és a
dir, per als quals es compleix la relació\[
f(x)=g(x)
\]En el nostre cas, tenim\[
\begin{array}{ccc}
-x+1=2x-4 & \Longrightarrow & x=\frac{5}{3}\end{array}\]i, a més, la imatge de $5/3$ és\[
f(5/3)=-\frac{5}{3}+1=-\frac{2}{3}
\]Aleshores, les dues gràfiques es tallen en el punt de coordenades $(5/3,-2/3)$.

Com que es tracta de funcions afins, la representació gràfica de les quals són rectes, esbrinarem els valors de $x$ per als quals la funció és
positiva o negativa a partir del valor en el que la funció s'anul·la
(talla l'eix $X$). A partir d'aquest valor, segons sigui la funció
creixent o decreixent, determinarem els valors de $x$ per als quals la
funció és positiva i negativa. Així, la funció $f(x)=-x+1$ és
decreixent, perquè la seva gràfica és una recta de pendent negativa, i com que
s'anul·la en $x=1$, aleshores $f$ és positiva per a $x<1$ i negativa per a $x>1$. En canvi, la funció $g(x)=2x-4$ és creixent, perquè la seva gràfica és
una recta de pendent positiva, i com que s'anul·la en $x=2$, aleshores $g$ és
positiva per a $x>2$ i negativa per a $x<2$.
\end{solucio}

\begin{exercici}
Construeix la gràfica de les funcions les imatges de les quals estan
definides per les relacions següents: (a) $y=x^{2}-3$; (b) $y=-x^{2}+5$;
(c) $y=(x+2)^{2}-3$; (d) $y=-2(x-3)^{2}+3$. Esbrinar si tallen o no l'eix
d'abscisses.
\end{exercici}

\begin{solucio}
Evidentment, totes les funcions indicades són polinòmiques de segon
grau i, per tant, tindran com a gràfica una paràbola.

a) La relació $y=x^{2}-3$ estableix una funció quadràtica $f$
definida per $f(x)=x^{2}-3$, per a tot $x\in \mathbb{R}$. La gràfica d'
aquesta funció s'obtindrà a partir de la paràbola $y=x^{2}$
mitjançant una translació seguint l'eix d’ordenades $3$ unitats en el
sentit negatiu. \begin{center}
\includegraphics[width=2.6109in,height=2.559in,keepaspectratio]{poli17.png}
\end{center}A partir de la gràfica observem que hi ha dos talls amb el
eix d’abscisses. Aquests punts són els zeros de la funció, és a dir, els
valors de $x$ que verifiquen l'equació\[
x^{2}-3=0
\]les quals solucions són $\pm \sqrt{3}$.

b) La relació $y=-x^{2}+5$ estableix una funció quadràtica $f$
definida per $f(x)=-x^{2}+5$, per a tot $x\in \mathbb{R}$. La gràfica d'
aquesta funció s'obtindrà a partir de la paràbola $y=-x^{2}$
mitjançant una translació seguint l'eix d’ordenades $5$ unitats en el
sentit positiu.\begin{center}
\includegraphics[width=2.5486in,height=2.5382in,keepaspectratio]{poli18.png}
\end{center}A partir de la gràfica observem que hi ha dos talls amb el
eix d’abscisses. Aquests punts són els zeros de la funció, és a dir, els
valors de $x$ que verifiquen l'equació\[
-x^{2}+5=0
\]les quals solucions són $\pm \sqrt{5}$.

c) La relació $y=(x+2)^{2}-3$ estableix una funció quadràtica $f$ definida per $f(x)=(x+2)^{2}-3$, per a tot $x\in \mathbb{R}$. La gràfica d'aquesta funció s'obtindrà a partir de la paràbola $y=x^{2} $fent primer una translació seguint l'eix d’ordenades $3$ unitats en el sentit negatiu, i, després, una translació seguint
l'eix d'abscisses $2$ unitats en el sentit negatiu. \begin{center}
\includegraphics[width=2.9853in,height=2.5382in,keepaspectratio]{poli19.png}
\end{center}A partir
de la gràfica observem que hi ha dos talls amb l'eix d’abscisses.
Aquests punts són els zeros de la funció, és a dir, els valors de $x$
que verifiquen l'equació\begin{eqnarray*}
(x+2)^{2}-3 &=&0 \\
(x+2)^{2} &=&3
\end{eqnarray*}les quals solucions són $-2\pm \sqrt{3}$.

d) La relació $y=-2(x-3)^{2}+3$ estableix una funció quadràtica $f$ definida per $f(x)=-2(x-3)^{2}+3$, per a tot $x\in \mathbb{R}$. La
gràfica d'aquesta funció s'obtindrà a partir de la paràbola $y=-2x^{2}$ fent primer una translació seguint l'eix d’ordenades
$3$ unitats en el sentit positiu, i, després, una translació seguint
l'eix d'abscisses $3$ unitats en el sentit positiu. \begin{center}
\includegraphics[width=2.8297in,height=2.6948in,keepaspectratio]{poli20.png}
\end{center}A partir
de la gràfica observem que hi ha dos talls amb l'eix d’abscisses.
Aquests punts són els zeros de la funció, és a dir, els valors de $x$
que verifiquen l'equació\begin{eqnarray*}
-2(x-3)^{2}+3 &=&0 \\
2(x-3)^{2} &=&3 \\
(x-3)^{2} &=&\frac{3}{2}
\end{eqnarray*}les quals solucions són $3\pm \sqrt{3/2}$.
\end{solucio}

\begin{exercici}
Representa les gràfiques de les funcions següents, les imatges de les quals
estan definides per les relacions següents: (a) $y=x^{2}-2x-3$; (b) $y=-2x^{2}+4x-3$; i (c) $y=2x^{2}-4x$. Caracteritza cada una d'aquestes gràfiques.
\end{exercici}

\begin{solucio}
Totes les funcions indicades són quadràtiques i, per tant, tindran
com a gràfica una paràbola.

a) La relació $y=x^{2}-2x-3$ estableix una funció quadràtica $f $definida per $f(x)=x^{2}-2x-3$, per a tot $x\in \mathbb{R}$. Observa que
podem escriure\[
x^{2}-2x-3=(x-1)^{2}-4
\]i, per tant, la gràfica d'aquesta funció s'obtindrà a partir de
la paràbola $y=x^{2}$ fent primer una translació de $4$
unitats en la direcció de l'eix $Y$ en sentit negatiu, i, després, una
translació de $1$ unitat en la direcció de l'eix $X$ en sentit
positiu. \begin{center}
\includegraphics[width=2.642in,height=2.5486in,keepaspectratio]{poli21.png}
\end{center}És clar que el vèrtex de la paràbola es troba en el
punt de coordenades $(1,-4)$, que és el transformat de $(0,0)$ per les
operacions indicades abans. L'eix de simetria és la recta vertical $x=1$. La funció té un punt d'ordenada mínima en $x=1$, el valor del qual és $-4$. La funció és creixent per a $x>1$ i decreixent per a $x<1 $. La gràfica talla l'eix $Y$ en el punt $(0,-3)$ i l'eix $X$ en
els punts d'abscissa que verifiquen l'equació\[
x^{2}-2x-3=0
\]és a dir, en $x=-1$ i en $x=$ $3$. La funció és positiva en $(-\infty,-1)\cup (3,+\infty )$ i negativa en $(-1,3)$.

b) La relació $y=-2x^{2}+4x-3$ estableix una funció quadràtica
$f$ definida per $f(x)=-2x^{2}+4x-3$, per a tot $x\in \mathbb{R}$. Observa
que podem escriure\begin{eqnarray*}
-2x^{2}+4x-3 &=&-2\left( x^{2}-2x+\frac{3}{2}\right) \\
&=&-2\left[ \left( x-1\right) ^{2}+\frac{1}{2}\right] \\
&=&-2(x-1)^{2}-1
\end{eqnarray*}i, per tant, la gràfica d'aquesta funció s'obtindrà a partir de
la paràbola $y=-2x^{2}$ fent primer una translació de $1$
unitat en la direcció de l'eix $Y$ en sentit negatiu, i, després, una
translació de $1$ unitat en la direcció de l'eix $X$ en sentit
positiu. \begin{center}
\includegraphics[width=2.4751in,height=2.674in,keepaspectratio]{poli22.png}
\end{center}És clar que el vèrtex de la paràbola es troba en el
punt de coordenades $(1,-1)$, que és el transformat de $(0,0)$ per les
operacions abans indicades. L'eix de simetria és la recta vertical $x=1$. La funció té un punt d'ordenada màxima en $x=1$, el valor del qual és $-1$. La funció és creixent per a $x<1$ i decreixent per a $x>1 $. La gràfica talla l'eix $Y$ en el punt $(0,-3)$ i no talla l’
eix $X$. La funció és sempre negativa.

c) La relació $y=2x^{2}-4x$ estableix una funció quadràtica $f$
definida per $f(x)=2x^{2}-4x$, per a tot $x\in \mathbb{R}$. Observa que
podem escriure\begin{eqnarray*}
2x^{2}-4x &=&2(x^{2}-2x) \\
&=&2\left[ (x-1)^{2}-1\right] \\
&=&2(x-1)^{2}-2
\end{eqnarray*}i, per tant, la gràfica d'aquesta funció s'obtindrà a partir de
la paràbola $y=2x^{2}$ fent primer una translació de $2$
unitats en la direcció de l'eix $Y$ en sentit negatiu, i, després, una
translació de $1$ unitat en la direcció de l'eix $X$ en sentit
positiu. \begin{center}
\includegraphics[width=2.559in,height=2.5694in,keepaspectratio]{poli23.png}
\end{center}És clar que el vèrtex de la paràbola es troba en el
punt de coordenades $(1,-2)$, que és el transformat de $(0,0)$ per les
operacions indicades abans. L'eix de simetria és la recta vertical $x=1$. La funció té un punt d'ordenada mínima en $x=1$, el valor del qual és $-2$. La funció és creixent per a $x>1$ i decreixent per a $x<1 $. La gràfica talla l'eix $Y$ en el punt $(0,0)$ i l'eix $X$ en
els punts d'abscissa que verifiquen l'equació\begin{eqnarray*}
2x^{2}-4x &=&0 \\
2x(x-2) &=&0
\end{eqnarray*}és a dir, en $x=0$ i en $x=$ $2$. La funció és positiva en $(-\infty,0)\cup (2,+\infty )$ i negativa en $(0,2)$.
\end{solucio}

\begin{exercici}
Construeix les gràfiques de les funcions següents: (a) $f(x)=x^{3}-5$;
(b) $f(x)=-x^{4}+5$; i (c) $f(x)=(x+3)^{3}+2$.
\end{exercici}

\begin{solucio}
Totes les funcions d'aquest exercici són polinòmiques de grau més gran que
$2$.

a) La gràfica de la funció $f(x)=x^{3}-5$ s'obtindrà a partir
de la de $y=x^{3}$ mitjançant una translació de $5$ unitats en la direcció de l'eix d’ordenades en sentit negatiu.\begin{center}
\includegraphics[width=2.4751in,height=2.642in,keepaspectratio]{poli24.png}
\end{center}

b) La gràfica de $f(x)=-x^{4}+5$ s'obtindrà a partir de la de $y=x^{4}$, fent primer una simetria respecte de l'eix $X$ i després traslladant-la $5$ unitats en la direcció de l'eix d’ordenades en
el sentit positiu. \begin{center}
\includegraphics[width=2.5278in,height=2.6109in,keepaspectratio]{poli25.png}
\end{center}

c) La gràfica de $f(x)=(x+3)^{3}+2$ s'obtindrà a partir de la de $y=x^{3}$ fent primer una translació de $2$ unitats en la direcció de l'eix $Y$ en sentit positiu i, després, una translació de
$3$ unitats en la direcció de l'eix $X$ en el sentit negatiu.\begin{center}
\includegraphics[width=2.5694in,height=2.6005in,keepaspectratio]{poli26.png}
\end{center}
\end{solucio}
